\documentclass[11pt]{article}

    \usepackage[breakable]{tcolorbox}
    \usepackage{parskip} % Stop auto-indenting (to mimic markdown behaviour)
    

    % Basic figure setup, for now with no caption control since it's done
    % automatically by Pandoc (which extracts ![](path) syntax from Markdown).
    \usepackage{graphicx}
    % Keep aspect ratio if custom image width or height is specified
    \setkeys{Gin}{keepaspectratio}
    % Maintain compatibility with old templates. Remove in nbconvert 6.0
    \let\Oldincludegraphics\includegraphics
    % Ensure that by default, figures have no caption (until we provide a
    % proper Figure object with a Caption API and a way to capture that
    % in the conversion process - todo).
    \usepackage{caption}
    \DeclareCaptionFormat{nocaption}{}
    \captionsetup{format=nocaption,aboveskip=0pt,belowskip=0pt}

    \usepackage{float}
    \floatplacement{figure}{H} % forces figures to be placed at the correct location
    \usepackage{xcolor} % Allow colors to be defined
    \usepackage{enumerate} % Needed for markdown enumerations to work
    \usepackage{geometry} % Used to adjust the document margins
    \usepackage{amsmath} % Equations
    \usepackage{amssymb} % Equations
    \usepackage{textcomp} % defines textquotesingle
    % Hack from http://tex.stackexchange.com/a/47451/13684:
    \AtBeginDocument{%
        \def\PYZsq{\textquotesingle}% Upright quotes in Pygmentized code
    }
    \usepackage{upquote} % Upright quotes for verbatim code
    \usepackage{eurosym} % defines \euro

    \usepackage{iftex}
    \ifPDFTeX
        \usepackage[T1]{fontenc}
        \IfFileExists{alphabeta.sty}{
              \usepackage{alphabeta}
          }{
              \usepackage[mathletters]{ucs}
              \usepackage[utf8x]{inputenc}
          }
    \else
        \usepackage{fontspec}
        \usepackage{unicode-math}
    \fi

    \usepackage{fancyvrb} % verbatim replacement that allows latex
    \usepackage{grffile} % extends the file name processing of package graphics
                         % to support a larger range
    \makeatletter % fix for old versions of grffile with XeLaTeX
    \@ifpackagelater{grffile}{2019/11/01}
    {
      % Do nothing on new versions
    }
    {
      \def\Gread@@xetex#1{%
        \IfFileExists{"\Gin@base".bb}%
        {\Gread@eps{\Gin@base.bb}}%
        {\Gread@@xetex@aux#1}%
      }
    }
    \makeatother
    \usepackage[Export]{adjustbox} % Used to constrain images to a maximum size
    \adjustboxset{max size={0.9\linewidth}{0.9\paperheight}}

    % The hyperref package gives us a pdf with properly built
    % internal navigation ('pdf bookmarks' for the table of contents,
    % internal cross-reference links, web links for URLs, etc.)
    \usepackage{hyperref}
    % The default LaTeX title has an obnoxious amount of whitespace. By default,
    % titling removes some of it. It also provides customization options.
    \usepackage{titling}
    \usepackage{longtable} % longtable support required by pandoc >1.10
    \usepackage{booktabs}  % table support for pandoc > 1.12.2
    \usepackage{array}     % table support for pandoc >= 2.11.3
    \usepackage{calc}      % table minipage width calculation for pandoc >= 2.11.1
    \usepackage[inline]{enumitem} % IRkernel/repr support (it uses the enumerate* environment)
    \usepackage[normalem]{ulem} % ulem is needed to support strikethroughs (\sout)
                                % normalem makes italics be italics, not underlines
    \usepackage{soul}      % strikethrough (\st) support for pandoc >= 3.0.0
    \usepackage{mathrsfs}
    

    
    % Colors for the hyperref package
    \definecolor{urlcolor}{rgb}{0,.145,.698}
    \definecolor{linkcolor}{rgb}{.71,0.21,0.01}
    \definecolor{citecolor}{rgb}{.12,.54,.11}

    % ANSI colors
    \definecolor{ansi-black}{HTML}{3E424D}
    \definecolor{ansi-black-intense}{HTML}{282C36}
    \definecolor{ansi-red}{HTML}{E75C58}
    \definecolor{ansi-red-intense}{HTML}{B22B31}
    \definecolor{ansi-green}{HTML}{00A250}
    \definecolor{ansi-green-intense}{HTML}{007427}
    \definecolor{ansi-yellow}{HTML}{DDB62B}
    \definecolor{ansi-yellow-intense}{HTML}{B27D12}
    \definecolor{ansi-blue}{HTML}{208FFB}
    \definecolor{ansi-blue-intense}{HTML}{0065CA}
    \definecolor{ansi-magenta}{HTML}{D160C4}
    \definecolor{ansi-magenta-intense}{HTML}{A03196}
    \definecolor{ansi-cyan}{HTML}{60C6C8}
    \definecolor{ansi-cyan-intense}{HTML}{258F8F}
    \definecolor{ansi-white}{HTML}{C5C1B4}
    \definecolor{ansi-white-intense}{HTML}{A1A6B2}
    \definecolor{ansi-default-inverse-fg}{HTML}{FFFFFF}
    \definecolor{ansi-default-inverse-bg}{HTML}{000000}

    % common color for the border for error outputs.
    \definecolor{outerrorbackground}{HTML}{FFDFDF}

    % commands and environments needed by pandoc snippets
    % extracted from the output of `pandoc -s`
    \providecommand{\tightlist}{%
      \setlength{\itemsep}{0pt}\setlength{\parskip}{0pt}}
    \DefineVerbatimEnvironment{Highlighting}{Verbatim}{commandchars=\\\{\}}
    % Add ',fontsize=\small' for more characters per line
    \newenvironment{Shaded}{}{}
    \newcommand{\KeywordTok}[1]{\textcolor[rgb]{0.00,0.44,0.13}{\textbf{{#1}}}}
    \newcommand{\DataTypeTok}[1]{\textcolor[rgb]{0.56,0.13,0.00}{{#1}}}
    \newcommand{\DecValTok}[1]{\textcolor[rgb]{0.25,0.63,0.44}{{#1}}}
    \newcommand{\BaseNTok}[1]{\textcolor[rgb]{0.25,0.63,0.44}{{#1}}}
    \newcommand{\FloatTok}[1]{\textcolor[rgb]{0.25,0.63,0.44}{{#1}}}
    \newcommand{\CharTok}[1]{\textcolor[rgb]{0.25,0.44,0.63}{{#1}}}
    \newcommand{\StringTok}[1]{\textcolor[rgb]{0.25,0.44,0.63}{{#1}}}
    \newcommand{\CommentTok}[1]{\textcolor[rgb]{0.38,0.63,0.69}{\textit{{#1}}}}
    \newcommand{\OtherTok}[1]{\textcolor[rgb]{0.00,0.44,0.13}{{#1}}}
    \newcommand{\AlertTok}[1]{\textcolor[rgb]{1.00,0.00,0.00}{\textbf{{#1}}}}
    \newcommand{\FunctionTok}[1]{\textcolor[rgb]{0.02,0.16,0.49}{{#1}}}
    \newcommand{\RegionMarkerTok}[1]{{#1}}
    \newcommand{\ErrorTok}[1]{\textcolor[rgb]{1.00,0.00,0.00}{\textbf{{#1}}}}
    \newcommand{\NormalTok}[1]{{#1}}

    % Additional commands for more recent versions of Pandoc
    \newcommand{\ConstantTok}[1]{\textcolor[rgb]{0.53,0.00,0.00}{{#1}}}
    \newcommand{\SpecialCharTok}[1]{\textcolor[rgb]{0.25,0.44,0.63}{{#1}}}
    \newcommand{\VerbatimStringTok}[1]{\textcolor[rgb]{0.25,0.44,0.63}{{#1}}}
    \newcommand{\SpecialStringTok}[1]{\textcolor[rgb]{0.73,0.40,0.53}{{#1}}}
    \newcommand{\ImportTok}[1]{{#1}}
    \newcommand{\DocumentationTok}[1]{\textcolor[rgb]{0.73,0.13,0.13}{\textit{{#1}}}}
    \newcommand{\AnnotationTok}[1]{\textcolor[rgb]{0.38,0.63,0.69}{\textbf{\textit{{#1}}}}}
    \newcommand{\CommentVarTok}[1]{\textcolor[rgb]{0.38,0.63,0.69}{\textbf{\textit{{#1}}}}}
    \newcommand{\VariableTok}[1]{\textcolor[rgb]{0.10,0.09,0.49}{{#1}}}
    \newcommand{\ControlFlowTok}[1]{\textcolor[rgb]{0.00,0.44,0.13}{\textbf{{#1}}}}
    \newcommand{\OperatorTok}[1]{\textcolor[rgb]{0.40,0.40,0.40}{{#1}}}
    \newcommand{\BuiltInTok}[1]{{#1}}
    \newcommand{\ExtensionTok}[1]{{#1}}
    \newcommand{\PreprocessorTok}[1]{\textcolor[rgb]{0.74,0.48,0.00}{{#1}}}
    \newcommand{\AttributeTok}[1]{\textcolor[rgb]{0.49,0.56,0.16}{{#1}}}
    \newcommand{\InformationTok}[1]{\textcolor[rgb]{0.38,0.63,0.69}{\textbf{\textit{{#1}}}}}
    \newcommand{\WarningTok}[1]{\textcolor[rgb]{0.38,0.63,0.69}{\textbf{\textit{{#1}}}}}
    \makeatletter
    \newsavebox\pandoc@box
    \newcommand*\pandocbounded[1]{%
      \sbox\pandoc@box{#1}%
      % scaling factors for width and height
      \Gscale@div\@tempa\textheight{\dimexpr\ht\pandoc@box+\dp\pandoc@box\relax}%
      \Gscale@div\@tempb\linewidth{\wd\pandoc@box}%
      % select the smaller of both
      \ifdim\@tempb\p@<\@tempa\p@
        \let\@tempa\@tempb
      \fi
      % scaling accordingly (\@tempa < 1)
      \ifdim\@tempa\p@<\p@
        \scalebox{\@tempa}{\usebox\pandoc@box}%
      % scaling not needed, use as it is
      \else
        \usebox{\pandoc@box}%
      \fi
    }
    \makeatother

    % Define a nice break command that doesn't care if a line doesn't already
    % exist.
    \def\br{\hspace*{\fill} \\* }
    % Math Jax compatibility definitions
    \def\gt{>}
    \def\lt{<}
    \let\Oldtex\TeX
    \let\Oldlatex\LaTeX
    \renewcommand{\TeX}{\textrm{\Oldtex}}
    \renewcommand{\LaTeX}{\textrm{\Oldlatex}}
    % Document parameters
    % Document title
    \title{dagamesolve: A Package for Finding All Nash Equilibria in Sequences of Discrete-Action Games in Stata}
    
    
    \date{Draft: October 2026}
    
    
    
    
    \author{Matthew J. Baker\\[0.25em]\small Hunter College and the Graduate Center, CUNY}
    
    
    
% Pygments definitions
\makeatletter
\def\PY@reset{\let\PY@it=\relax \let\PY@bf=\relax%
    \let\PY@ul=\relax \let\PY@tc=\relax%
    \let\PY@bc=\relax \let\PY@ff=\relax}
\def\PY@tok#1{\csname PY@tok@#1\endcsname}
\def\PY@toks#1+{\ifx\relax#1\empty\else%
    \PY@tok{#1}\expandafter\PY@toks\fi}
\def\PY@do#1{\PY@bc{\PY@tc{\PY@ul{%
    \PY@it{\PY@bf{\PY@ff{#1}}}}}}}
\def\PY#1#2{\PY@reset\PY@toks#1+\relax+\PY@do{#2}}

\@namedef{PY@tok@w}{\def\PY@tc##1{\textcolor[rgb]{0.73,0.73,0.73}{##1}}}
\@namedef{PY@tok@c}{\let\PY@it=\textit\def\PY@tc##1{\textcolor[rgb]{0.24,0.48,0.48}{##1}}}
\@namedef{PY@tok@cp}{\def\PY@tc##1{\textcolor[rgb]{0.61,0.40,0.00}{##1}}}
\@namedef{PY@tok@k}{\let\PY@bf=\textbf\def\PY@tc##1{\textcolor[rgb]{0.00,0.50,0.00}{##1}}}
\@namedef{PY@tok@kp}{\def\PY@tc##1{\textcolor[rgb]{0.00,0.50,0.00}{##1}}}
\@namedef{PY@tok@kt}{\def\PY@tc##1{\textcolor[rgb]{0.69,0.00,0.25}{##1}}}
\@namedef{PY@tok@o}{\def\PY@tc##1{\textcolor[rgb]{0.40,0.40,0.40}{##1}}}
\@namedef{PY@tok@ow}{\let\PY@bf=\textbf\def\PY@tc##1{\textcolor[rgb]{0.67,0.13,1.00}{##1}}}
\@namedef{PY@tok@nb}{\def\PY@tc##1{\textcolor[rgb]{0.00,0.50,0.00}{##1}}}
\@namedef{PY@tok@nf}{\def\PY@tc##1{\textcolor[rgb]{0.00,0.00,1.00}{##1}}}
\@namedef{PY@tok@nc}{\let\PY@bf=\textbf\def\PY@tc##1{\textcolor[rgb]{0.00,0.00,1.00}{##1}}}
\@namedef{PY@tok@nn}{\let\PY@bf=\textbf\def\PY@tc##1{\textcolor[rgb]{0.00,0.00,1.00}{##1}}}
\@namedef{PY@tok@ne}{\let\PY@bf=\textbf\def\PY@tc##1{\textcolor[rgb]{0.80,0.25,0.22}{##1}}}
\@namedef{PY@tok@nv}{\def\PY@tc##1{\textcolor[rgb]{0.10,0.09,0.49}{##1}}}
\@namedef{PY@tok@no}{\def\PY@tc##1{\textcolor[rgb]{0.53,0.00,0.00}{##1}}}
\@namedef{PY@tok@nl}{\def\PY@tc##1{\textcolor[rgb]{0.46,0.46,0.00}{##1}}}
\@namedef{PY@tok@ni}{\let\PY@bf=\textbf\def\PY@tc##1{\textcolor[rgb]{0.44,0.44,0.44}{##1}}}
\@namedef{PY@tok@na}{\def\PY@tc##1{\textcolor[rgb]{0.41,0.47,0.13}{##1}}}
\@namedef{PY@tok@nt}{\let\PY@bf=\textbf\def\PY@tc##1{\textcolor[rgb]{0.00,0.50,0.00}{##1}}}
\@namedef{PY@tok@nd}{\def\PY@tc##1{\textcolor[rgb]{0.67,0.13,1.00}{##1}}}
\@namedef{PY@tok@s}{\def\PY@tc##1{\textcolor[rgb]{0.73,0.13,0.13}{##1}}}
\@namedef{PY@tok@sd}{\let\PY@it=\textit\def\PY@tc##1{\textcolor[rgb]{0.73,0.13,0.13}{##1}}}
\@namedef{PY@tok@si}{\let\PY@bf=\textbf\def\PY@tc##1{\textcolor[rgb]{0.64,0.35,0.47}{##1}}}
\@namedef{PY@tok@se}{\let\PY@bf=\textbf\def\PY@tc##1{\textcolor[rgb]{0.67,0.36,0.12}{##1}}}
\@namedef{PY@tok@sr}{\def\PY@tc##1{\textcolor[rgb]{0.64,0.35,0.47}{##1}}}
\@namedef{PY@tok@ss}{\def\PY@tc##1{\textcolor[rgb]{0.10,0.09,0.49}{##1}}}
\@namedef{PY@tok@sx}{\def\PY@tc##1{\textcolor[rgb]{0.00,0.50,0.00}{##1}}}
\@namedef{PY@tok@m}{\def\PY@tc##1{\textcolor[rgb]{0.40,0.40,0.40}{##1}}}
\@namedef{PY@tok@gh}{\let\PY@bf=\textbf\def\PY@tc##1{\textcolor[rgb]{0.00,0.00,0.50}{##1}}}
\@namedef{PY@tok@gu}{\let\PY@bf=\textbf\def\PY@tc##1{\textcolor[rgb]{0.50,0.00,0.50}{##1}}}
\@namedef{PY@tok@gd}{\def\PY@tc##1{\textcolor[rgb]{0.63,0.00,0.00}{##1}}}
\@namedef{PY@tok@gi}{\def\PY@tc##1{\textcolor[rgb]{0.00,0.52,0.00}{##1}}}
\@namedef{PY@tok@gr}{\def\PY@tc##1{\textcolor[rgb]{0.89,0.00,0.00}{##1}}}
\@namedef{PY@tok@ge}{\let\PY@it=\textit}
\@namedef{PY@tok@gs}{\let\PY@bf=\textbf}
\@namedef{PY@tok@gp}{\let\PY@bf=\textbf\def\PY@tc##1{\textcolor[rgb]{0.00,0.00,0.50}{##1}}}
\@namedef{PY@tok@go}{\def\PY@tc##1{\textcolor[rgb]{0.44,0.44,0.44}{##1}}}
\@namedef{PY@tok@gt}{\def\PY@tc##1{\textcolor[rgb]{0.00,0.27,0.87}{##1}}}
\@namedef{PY@tok@err}{\def\PY@bc##1{{\setlength{\fboxsep}{\string -\fboxrule}\fcolorbox[rgb]{1.00,0.00,0.00}{1,1,1}{\strut ##1}}}}
\@namedef{PY@tok@kc}{\let\PY@bf=\textbf\def\PY@tc##1{\textcolor[rgb]{0.00,0.50,0.00}{##1}}}
\@namedef{PY@tok@kd}{\let\PY@bf=\textbf\def\PY@tc##1{\textcolor[rgb]{0.00,0.50,0.00}{##1}}}
\@namedef{PY@tok@kn}{\let\PY@bf=\textbf\def\PY@tc##1{\textcolor[rgb]{0.00,0.50,0.00}{##1}}}
\@namedef{PY@tok@kr}{\let\PY@bf=\textbf\def\PY@tc##1{\textcolor[rgb]{0.00,0.50,0.00}{##1}}}
\@namedef{PY@tok@bp}{\def\PY@tc##1{\textcolor[rgb]{0.00,0.50,0.00}{##1}}}
\@namedef{PY@tok@fm}{\def\PY@tc##1{\textcolor[rgb]{0.00,0.00,1.00}{##1}}}
\@namedef{PY@tok@vc}{\def\PY@tc##1{\textcolor[rgb]{0.10,0.09,0.49}{##1}}}
\@namedef{PY@tok@vg}{\def\PY@tc##1{\textcolor[rgb]{0.10,0.09,0.49}{##1}}}
\@namedef{PY@tok@vi}{\def\PY@tc##1{\textcolor[rgb]{0.10,0.09,0.49}{##1}}}
\@namedef{PY@tok@vm}{\def\PY@tc##1{\textcolor[rgb]{0.10,0.09,0.49}{##1}}}
\@namedef{PY@tok@sa}{\def\PY@tc##1{\textcolor[rgb]{0.73,0.13,0.13}{##1}}}
\@namedef{PY@tok@sb}{\def\PY@tc##1{\textcolor[rgb]{0.73,0.13,0.13}{##1}}}
\@namedef{PY@tok@sc}{\def\PY@tc##1{\textcolor[rgb]{0.73,0.13,0.13}{##1}}}
\@namedef{PY@tok@dl}{\def\PY@tc##1{\textcolor[rgb]{0.73,0.13,0.13}{##1}}}
\@namedef{PY@tok@s2}{\def\PY@tc##1{\textcolor[rgb]{0.73,0.13,0.13}{##1}}}
\@namedef{PY@tok@sh}{\def\PY@tc##1{\textcolor[rgb]{0.73,0.13,0.13}{##1}}}
\@namedef{PY@tok@s1}{\def\PY@tc##1{\textcolor[rgb]{0.73,0.13,0.13}{##1}}}
\@namedef{PY@tok@mb}{\def\PY@tc##1{\textcolor[rgb]{0.40,0.40,0.40}{##1}}}
\@namedef{PY@tok@mf}{\def\PY@tc##1{\textcolor[rgb]{0.40,0.40,0.40}{##1}}}
\@namedef{PY@tok@mh}{\def\PY@tc##1{\textcolor[rgb]{0.40,0.40,0.40}{##1}}}
\@namedef{PY@tok@mi}{\def\PY@tc##1{\textcolor[rgb]{0.40,0.40,0.40}{##1}}}
\@namedef{PY@tok@il}{\def\PY@tc##1{\textcolor[rgb]{0.40,0.40,0.40}{##1}}}
\@namedef{PY@tok@mo}{\def\PY@tc##1{\textcolor[rgb]{0.40,0.40,0.40}{##1}}}
\@namedef{PY@tok@ch}{\let\PY@it=\textit\def\PY@tc##1{\textcolor[rgb]{0.24,0.48,0.48}{##1}}}
\@namedef{PY@tok@cm}{\let\PY@it=\textit\def\PY@tc##1{\textcolor[rgb]{0.24,0.48,0.48}{##1}}}
\@namedef{PY@tok@cpf}{\let\PY@it=\textit\def\PY@tc##1{\textcolor[rgb]{0.24,0.48,0.48}{##1}}}
\@namedef{PY@tok@c1}{\let\PY@it=\textit\def\PY@tc##1{\textcolor[rgb]{0.24,0.48,0.48}{##1}}}
\@namedef{PY@tok@cs}{\let\PY@it=\textit\def\PY@tc##1{\textcolor[rgb]{0.24,0.48,0.48}{##1}}}

\def\PYZbs{\char`\\}
\def\PYZus{\char`\_}
\def\PYZob{\char`\{}
\def\PYZcb{\char`\}}
\def\PYZca{\char`\^}
\def\PYZam{\char`\&}
\def\PYZlt{\char`\<}
\def\PYZgt{\char`\>}
\def\PYZsh{\char`\#}
\def\PYZpc{\char`\%}
\def\PYZdl{\char`\$}
\def\PYZhy{\char`\-}
\def\PYZsq{\char`\'}
\def\PYZdq{\char`\"}
\def\PYZti{\char`\~}
% for compatibility with earlier versions
\def\PYZat{@}
\def\PYZlb{[}
\def\PYZrb{]}
\makeatother


    % For linebreaks inside Verbatim environment from package fancyvrb.
    \makeatletter
        \newbox\Wrappedcontinuationbox
        \newbox\Wrappedvisiblespacebox
        \newcommand*\Wrappedvisiblespace {\textcolor{red}{\textvisiblespace}}
        \newcommand*\Wrappedcontinuationsymbol {\textcolor{red}{\llap{\tiny$\m@th\hookrightarrow$}}}
        \newcommand*\Wrappedcontinuationindent {3ex }
        \newcommand*\Wrappedafterbreak {\kern\Wrappedcontinuationindent\copy\Wrappedcontinuationbox}
        % Take advantage of the already applied Pygments mark-up to insert
        % potential linebreaks for TeX processing.
        %        {, <, #, %, $, ' and ": go to next line.
        %        _, }, ^, &, >, - and ~: stay at end of broken line.
        % Use of \textquotesingle for straight quote.
        \newcommand*\Wrappedbreaksatspecials {%
            \def\PYGZus{\discretionary{\char`\_}{\Wrappedafterbreak}{\char`\_}}%
            \def\PYGZob{\discretionary{}{\Wrappedafterbreak\char`\{}{\char`\{}}%
            \def\PYGZcb{\discretionary{\char`\}}{\Wrappedafterbreak}{\char`\}}}%
            \def\PYGZca{\discretionary{\char`\^}{\Wrappedafterbreak}{\char`\^}}%
            \def\PYGZam{\discretionary{\char`\&}{\Wrappedafterbreak}{\char`\&}}%
            \def\PYGZlt{\discretionary{}{\Wrappedafterbreak\char`\<}{\char`\<}}%
            \def\PYGZgt{\discretionary{\char`\>}{\Wrappedafterbreak}{\char`\>}}%
            \def\PYGZsh{\discretionary{}{\Wrappedafterbreak\char`\#}{\char`\#}}%
            \def\PYGZpc{\discretionary{}{\Wrappedafterbreak\char`\%}{\char`\%}}%
            \def\PYGZdl{\discretionary{}{\Wrappedafterbreak\char`\$}{\char`\$}}%
            \def\PYGZhy{\discretionary{\char`\-}{\Wrappedafterbreak}{\char`\-}}%
            \def\PYGZsq{\discretionary{}{\Wrappedafterbreak\textquotesingle}{\textquotesingle}}%
            \def\PYGZdq{\discretionary{}{\Wrappedafterbreak\char`\"}{\char`\"}}%
            \def\PYGZti{\discretionary{\char`\~}{\Wrappedafterbreak}{\char`\~}}%
        }
        % Some characters . , ; ? ! / are not pygmentized.
        % This macro makes them "active" and they will insert potential linebreaks
        \newcommand*\Wrappedbreaksatpunct {%
            \lccode`\~`\.\lowercase{\def~}{\discretionary{\hbox{\char`\.}}{\Wrappedafterbreak}{\hbox{\char`\.}}}%
            \lccode`\~`\,\lowercase{\def~}{\discretionary{\hbox{\char`\,}}{\Wrappedafterbreak}{\hbox{\char`\,}}}%
            \lccode`\~`\;\lowercase{\def~}{\discretionary{\hbox{\char`\;}}{\Wrappedafterbreak}{\hbox{\char`\;}}}%
            \lccode`\~`\:\lowercase{\def~}{\discretionary{\hbox{\char`\:}}{\Wrappedafterbreak}{\hbox{\char`\:}}}%
            \lccode`\~`\?\lowercase{\def~}{\discretionary{\hbox{\char`\?}}{\Wrappedafterbreak}{\hbox{\char`\?}}}%
            \lccode`\~`\!\lowercase{\def~}{\discretionary{\hbox{\char`\!}}{\Wrappedafterbreak}{\hbox{\char`\!}}}%
            \lccode`\~`\/\lowercase{\def~}{\discretionary{\hbox{\char`\/}}{\Wrappedafterbreak}{\hbox{\char`\/}}}%
            \catcode`\.\active
            \catcode`\,\active
            \catcode`\;\active
            \catcode`\:\active
            \catcode`\?\active
            \catcode`\!\active
            \catcode`\/\active
            \lccode`\~`\~
        }
    \makeatother

    \let\OriginalVerbatim=\Verbatim
    \makeatletter
    \renewcommand{\Verbatim}[1][1]{%
        %\parskip\z@skip
        \sbox\Wrappedcontinuationbox {\Wrappedcontinuationsymbol}%
        \sbox\Wrappedvisiblespacebox {\FV@SetupFont\Wrappedvisiblespace}%
        \def\FancyVerbFormatLine ##1{\hsize\linewidth
            \vtop{\raggedright\hyphenpenalty\z@\exhyphenpenalty\z@
                \doublehyphendemerits\z@\finalhyphendemerits\z@
                \strut ##1\strut}%
        }%
        % If the linebreak is at a space, the latter will be displayed as visible
        % space at end of first line, and a continuation symbol starts next line.
        % Stretch/shrink are however usually zero for typewriter font.
        \def\FV@Space {%
            \nobreak\hskip\z@ plus\fontdimen3\font minus\fontdimen4\font
            \discretionary{\copy\Wrappedvisiblespacebox}{\Wrappedafterbreak}
            {\kern\fontdimen2\font}%
        }%

        % Allow breaks at special characters using \PYG... macros.
        \Wrappedbreaksatspecials
        % Breaks at punctuation characters . , ; ? ! and / need catcode=\active
        \OriginalVerbatim[#1,codes*=\Wrappedbreaksatpunct]%
    }
    \makeatother

    % Exact colors from NB
    \definecolor{incolor}{HTML}{303F9F}
    \definecolor{outcolor}{HTML}{D84315}
    \definecolor{cellborder}{HTML}{CFCFCF}
    \definecolor{cellbackground}{HTML}{F7F7F7}

    % prompt
    \makeatletter
    \newcommand{\boxspacing}{\kern\kvtcb@left@rule\kern\kvtcb@boxsep}
    \makeatother
    \newcommand{\prompt}[4]{
        {\ttfamily\llap{{\color{#2}[#3]:\hspace{3pt}#4}}\vspace{-\baselineskip}}
    }
    

    
    % Prevent overflowing lines due to hard-to-break entities
    \sloppy
    % Setup hyperref package
    \hypersetup{
      breaklinks=true,  % so long urls are correctly broken across lines
      colorlinks=true,
      urlcolor=urlcolor,
      linkcolor=linkcolor,
      citecolor=citecolor,
      }
    % Slightly bigger margins than the latex defaults
    
    \geometry{verbose,tmargin=1in,bmargin=1in,lmargin=1in,rmargin=1in}
    
    

\begin{document}
    
    \maketitle
    
    

    
    \textbf{Abstract.}\\
\texttt{dagamesolve} is a Stata/Mata package for solving finite,
simultaneous-move, discrete-action games of complete information. The
package supports an arbitrary number of players, allows players to have
different numbers of actions, searches for pure, partially mixed, and
fully mixed Nash equilibria, and can solve a sequence of games with a
common action structure in a single call.

The computational difficulty is not usually the detection of
pure-strategy equilibria. The hard part is the systematic search for
mixed-strategy equilibria. On a particular support, equilibrium requires
solving a system of nonlinear polynomial equations, and a complete
search must also consider the relevant supports. \texttt{dagamesolve}
combines dominance reduction, direct checks for pure equilibria, support
enumeration, interval methods for the nonlinear systems, and Newton
refinement. A faster point-based Newton search is also available.

This document explains the representation of games used by the package,
describes the numerical ideas underlying the solver, and provides
executable Stata examples. In particular, it emphasizes a feature that
is useful for econometric applications: a sequence of related games can
be constructed and solved directly in a Stata dataset.

    \begin{quote}
\textbf{About this notebook.}\\
This is intended to be both a technical note and an executable document.
Markdown cells contain the exposition and equations; Stata cells contain
examples that can be run directly with a Stata Jupyter kernel. The
notebook is the master source. A PDF can later be generated from the
executed notebook.

The theoretical discussion is adapted from the earlier presentation
\emph{Finding all solutions to discrete games (using Stata)} and updated
to match the current \texttt{dagamesolve} and \texttt{dagamestrats}
interfaces.
\end{quote}

\begin{quote}
To run the examples, install the package and its dependencies first; see
\textbf{Appendix A} for instructions and software checks. The notebook
initializes Stata just before the first example.
\end{quote}

    \section{Why solve finite games inside
Stata?}\label{why-solve-finite-games-inside-stata}

There are at least three motivations.

First, finding all Nash equilibria of a finite game is a nontrivial
computational problem, to put it mildly. Pure-strategy profiles can be
checked directly, but mixed strategies require solving nonlinear systems
and considering alternative supports. More generally, equilibrium
computation is known to be computationally difficult. Nash's theorem
guarantees that every finite game possesses at least one mixed-strategy
equilibrium, so the central computational problem is not one of deciding
whether a solution exists, but of actually finding one. Computing even a
single Nash equilibrium of a general finite game is PPAD-complete, a
complexity class associated with Brouwer fixed-point problems
(Daskalakis, Goldberg, and Papadimitriou, 2009). The problem of
recovering the entire equilibrium correspondence is potentially more
demanding still: even counting the Nash equilibria of a normal-form game
is \#P-hard (Conitzer and Sandholm, 2008).

Second, many empirical models require finding the equilibria for a
sequence of games, not just for a single isolated game. Entry models are
a natural example. Different markets may contain different payoff
shifters, while the underlying strategic structure and structural
parameters are common across markets. This repeated-game structure is
central to a large empirical literature on static discrete games; see
Aradillas-López (2020) for a recent review. Applications include models
of entry and market structure in retail, airlines, motels, video rental,
and other geographically differentiated industries (Bresnahan and Reiss,
1990, 1991; Berry, 1992; Mazzeo, 2002; Seim, 2006; Jia, 2008; Ciliberto
and Tamer, 2009). Bajari, Hong, and Ryan (2010), in particular, develop
an estimator for complete-information discrete games in which Nash
equilibria must be computed repeatedly as the structural parameters
vary. If the game solver is embedded in Stata, game construction,
equilibrium computation, data analysis, and eventually estimation can
all live in the same environment.

Third, sequences of games are useful and fun to study even outside
estimation. Part of my own interest in developing these routines came
from teaching graduate microeconomic theory at the CUNY Graduate Center.
Some familiar results in game theory become much more provocative when
one begins to ask what they imply quantitatively. Wilson's oddness
theorem (Wilson, 1971), for example, says that ``almost all'' finite
games have an odd number of Nash equilibria; the exceptional games form
a measure-zero set. The qualifier ``almost all'' naturally raises
further questions: how many equilibria does a typical game have, how is
that number distributed across games, and where do equilibria appear or
disappear as payoffs vary? Related questions arise in McLennan's (2005)
work on the expected number of Nash equilibria of a random normal-form
game. The magnitudes implied by those results can initially seem
surprising, particularly as the numbers of players or available actions
grow. Berg and McLennan (2005) pursue this question further for large
two-player normal-form games, deriving asymptotic results for the
expected number of equilibria. In any event, being able to generate and
solve large collections of games provides a simple computational
laboratory in which to explore such results directly: one can vary
payoff parameters, the number of players, or the number of available
actions and examine how equilibrium multiplicity changes.

These were the original motivations for developing \texttt{dagamesolve}:
not merely to solve one textbook game, but to make equilibrium
computation usable as a component inside a larger empirical workflow.
There was, however, an obvious reason not to develop yet another game
solver: \href{https://www.gambit-project.org/}{\textbf{Gambit}}, the
long-running and remarkably capable suite of software for the
construction and analysis of finite games developed by McKelvey,
McLennan, Turocy, and collaborators (2014). Gambit provides a rich
collection of equilibrium-computation algorithms, command-line tools, a
graphical interface, and a Python API, and it remains the natural
benchmark for general-purpose computation in finite games. My motivation
for writing a separate solver was therefore not that Gambit lacked the
relevant game-theoretic machinery. Rather, when this project began, I
wanted something that fit more naturally into the way I was already
doing empirical work in Stata: games represented directly as
observations and variables, sequences of related games solved in a
single call, and equilibrium results returned immediately to the working
dataset. I also wanted tighter control over the numerical treatment of
mixed-strategy solutions, particularly in \(N\)-player games where
support enumeration leads to systems of nonlinear polynomial equations.
Gambit itself offers exact rational arithmetic for several two-player
methods, while its general polynomial-equilibrium routines necessarily
rely on numerical procedures; \texttt{dagamesolve} took a different
route by combining support enumeration with interval arithmetic and
subsequent Newton refinement.

The last point of the previous paragraph leads me to a final point, and
an apology. \texttt{dagamesolve} is slow, in large part because of its
reliance on interval arithmetic. I will digress on the relative
advantages and disadvantages of this approach to solving polynomial
systems below, after first describing the nature of the problem.

    \section{Finite discrete-action
games}\label{finite-discrete-action-games}

Consider a complete-information simultaneous-move game with players
\(i \in N\) players where each player has a set of discrete actions:

\[
A_i=\{1,\ldots,K_i\}.
\]

A mixed strategy for player \(i\) is a probability vector

\[
\sigma_i=(\sigma_{i1},\ldots,\sigma_{iK_i}),
\qquad
\sigma_{ik}\geq 0,
\qquad
\sum_{k=1}^{K_i}\sigma_{ik}=1.
\]

A strategy profile \[\sigma^*=(\sigma_1^*,\ldots,\sigma_N^*)\] is a Nash
equilibrium if no player can increase expected payoff through a
unilateral deviation:

\[
u_i(\sigma_i^*,\sigma_{-i}^*)
\geq
u_i(\sigma_i',\sigma_{-i}^*)
\quad
\text{for every feasible }\sigma_i'.
\]

Pure strategies are special cases in which one action receives a weight
of probability equal to unity. A partially mixed equilibrium has some
players mixing and others playing pure strategies. A fully mixed
equilibrium places positive probability on every action in the relevant
support.

Mixed strategies matter computationally as well as theoretically. Nash's
existence theorem permits mixed strategies, and equilibrium search
cannot in general stop after inspecting the finite collection of pure
profiles.

    \subsection{A two-action illustration of
existence}\label{a-two-action-illustration-of-existence}

For intuition, suppose each player has two actions and write
\[\sigma_i\] for the probability of choosing action 1. One way of
motivating the fixed-point argument is to define the gain from shifting
toward action 1,

\[
g_i(\sigma)
=
\max\{u_i(1,\sigma_{-i})-u_i(2,\sigma_{-i}),0\},
\]

and construct a continuous adjustment mapping on the mixed-strategy
simplex. A fixed point of an appropriate mapping has the property that
profitable unilateral adjustments have been exhausted. This connection
between Nash equilibrium and fixed points underlies a number of
classical computational approaches to equilibrium.

An alternative formulation, due to McKelvey, turns the same idea into an
optimization problem. For a general finite game, let

\[
g_{ij}(\sigma)
=
\max\left\{
u_i(a_{ij},\sigma_{-i})-u_i(\sigma),0
\right\}
\]

denote the positive gain, if any, available to player (i) from deviating
to pure action (a\_\{ij\}). McKelvey constructs the nonnegative function

\[
V(\sigma) =\sum_i\sum_j[g_{ij}(\sigma)]^2.
\]

This function has the useful property that

\[
V(\sigma)=0
\quad\Longleftrightarrow\quad
\sigma\text{ is a Nash equilibrium}.
\]

Thus the computation of Nash equilibrium can alternatively be viewed as
a constrained global-minimization problem over the product of the
players' mixed-strategy simplices. McKelvey's formulation is
particularly attractive because the resulting objective is continuously
differentiable, allowing standard numerical optimization methods to be
brought to bear on the equilibrium problem (McKelvey, 1998; McKelvey and
McLennan, 1996)

The important computational lesson is not the particular fixed-point
proof. It is that equilibrium may lie in the interior of a probability
simplex rather than at one of its vertices. A numerical solver therefore
has to deal with continuous probability variables even though the
underlying actions are discrete.

The optimization formulation does not, however, make the problem of
finding all equilibria disappear. Numerical minimization begun from a
particular starting point can converge to one basin of attraction while
missing others, and a local minimum with a positive value of \(V\) is
not a Nash equilibrium. Searching for the entire equilibrium
correspondence therefore remains a global problem. This
distinction---between finding an equilibrium and systematically
searching for all equilibria---is central to the design of
\texttt{dagamesolve}.

    \section{Representing a game in ``list''
form}\label{representing-a-game-in-list-form}

Games are often displayed in normal form. For computation,
\texttt{dagamesolve} uses a list representation.

Let \(A\) denote a matrix whose columns enumerate the pure-strategy
action profiles and let \(\Pi\) contain the corresponding payoffs. If
there are \(N\) players, both matrices have \(N\) rows. The number of
columns is

\[
M=\prod_{i=1}^{N} K_i.
\]

For a two-player, two-action game, the action-profile matrix is

\[
A=
\begin{bmatrix}
1&1&2&2\\
1&2&1&2
\end{bmatrix}.
\]

Each column is one complete pure-strategy profile. The payoff matrix has
exactly the same column ordering.

In a Stata dataset, players are observations and action profiles are
variables. The utility command \texttt{dagamestrats} constructs the
action-profile variables automatically.

    \subsection{Example: constructing action
profiles}\label{example-constructing-action-profiles}

Suppose player 1 has two actions, player 2 has three, and player 3 has
two. There are \(2\times3\times2=12\) pure-strategy profiles.

    Before running the examples below, initialize the Stata session and set
the notebook graphics format. See Appendix A for installation
instructions.

    \begin{tcolorbox}[breakable, size=fbox, boxrule=1pt, pad at break*=1mm,colback=cellbackground, colframe=cellborder]
\prompt{In}{incolor}{72}{\boxspacing}
\begin{Verbatim}[commandchars=\\\{\}]
\PY{k}{clear} all
\PY{k}{set}\PY{k}{ more} off
\PY{k}{set} linesize \PY{l+m}{255}
\end{Verbatim}
\end{tcolorbox}

    \begin{tcolorbox}[breakable, size=fbox, boxrule=1pt, pad at break*=1mm,colback=cellbackground, colframe=cellborder]
\prompt{In}{incolor}{73}{\boxspacing}
\begin{Verbatim}[commandchars=\\\{\}]
\PYZpc{}set graph\PYZus{}format = png
\end{Verbatim}
\end{tcolorbox}

    \begin{Verbatim}[commandchars=\\\{\}]
graph\_format was png, is now png
    \end{Verbatim}

    \begin{tcolorbox}[breakable, size=fbox, boxrule=1pt, pad at break*=1mm,colback=cellbackground, colframe=cellborder]
\prompt{In}{incolor}{74}{\boxspacing}
\begin{Verbatim}[commandchars=\\\{\}]
\PYZpc{}set graph\PYZus{}width = 900px
\end{Verbatim}
\end{tcolorbox}

    \begin{tcolorbox}[breakable, size=fbox, boxrule=1pt, pad at break*=1mm,colback=cellbackground, colframe=cellborder]
\prompt{In}{incolor}{75}{\boxspacing}
\begin{Verbatim}[commandchars=\\\{\}]
\PY{k}{set} obs \PY{l+m}{3}

\PY{k}{gen} game = \PY{l+m}{1}
\PY{k}{gen} acts = \PY{l+m}{2}
\PY{k}{replace} acts = \PY{l+m}{3}\PY{k}{ in} \PY{l+m}{2}

dagamestrats acts, group(game)\PY{k}{ generate}(a)

\PY{k}{list} game acts a\PY{o}{*}
\PY{k}{return}\PY{k}{ list}
\end{Verbatim}
\end{tcolorbox}

    \begin{Verbatim}[commandchars=\\\{\}]
Number of observations (\_N) was 0, now 3.
(1 real change made)

     +--------------------------------------------------------------------------
---------+
     | game   acts   acts   a1   a2   a3   a4   a5   a6   a7   a8   a9   a10
a11   a12 |
     |--------------------------------------------------------------------------
---------|
  1. |    1      2      2    1    1    1    1    1    1    2    2    2     2
2     2 |
  2. |    1      3      3    1    1    2    2    3    3    1    1    2     2
3     3 |
  3. |    1      2      2    1    2    1    2    1    2    1    2    1     2
1     2 |
     +--------------------------------------------------------------------------
---------+

scalars:
           r(profiles) =  12
             r(groups) =  1
    \end{Verbatim}

    The generated variables \texttt{a1}, \texttt{a2}, \ldots, \texttt{a12}
are the twelve pure-strategy profiles. Corresponding payoff variables
can be created in the same order and passed to \texttt{dagamesolve}.

The same organization extends naturally to a \textbf{sequence of games}.
Each game occupies a group of observations, one observation per player.
Games in a single call must have the same number of players and the same
action structure, with corresponding players ordered consistently across
games.

    \section{Why mixed-strategy equilibrium is the difficult
part}\label{why-mixed-strategy-equilibrium-is-the-difficult-part}

Suppose, for simplicity, that three players each have two actions
\(a_i \in \{1,2\}\). Let \(p_i\) be the probability that player \(i\)
chooses action 1.

If all three players mix, each player must be indifferent between the
two actions. For player 1 this condition has the form

\[
\begin{aligned}
& p_2p_3\,u_1(1,1,1)
+p_2(1-p_3)\,u_1(1,1,2) \\
&\quad +(1-p_2)p_3\,u_1(1,2,1)
+(1-p_2)(1-p_3)\,u_1(1,2,2) \\
={}&
p_2p_3\,u_1(2,1,1)
+p_2(1-p_3)\,u_1(2,1,2) \\
&\quad +(1-p_2)p_3\,u_1(2,2,1)
+(1-p_2)(1-p_3)\,u_1(2,2,2).
\end{aligned}
\]

There is one analogous indifference equation for each mixing player.
Even in a binary-action game these equations are polynomial in the other
players' probabilities.

That is only part of the problem. A game may have equilibria on smaller
supports: one player may mix while another plays a pure strategy, or
different players may mix over different subsets of their actions. A
systematic search therefore has to consider candidate supports as well
as solve the nonlinear systems associated with them.

The number of pure profiles and the number of possible supports both
grow rapidly as players and actions are added.

    \section{\texorpdfstring{The solution strategy used by
\texttt{dagamesolve}}{The solution strategy used by dagamesolve}}\label{the-solution-strategy-used-by-dagamesolve}

At a high level, the current solver proceeds as follows.

\begin{enumerate}
\def\labelenumi{\arabic{enumi}.}
\item
  \textbf{Dominance reduction.} Iteratively eliminate strategies that
  can be removed by the package's dominance checks. In some games this
  step alone solves the game.
\item
  \textbf{Pure-strategy search.} Test the remaining pure profiles for
  Nash equilibrium.
\item
  \textbf{Support enumeration.} Construct candidate subgames/supports
  for mixed and partially mixed equilibria.
\item
  \textbf{Reduce each support.} Apply dominance checks again inside the
  candidate subgame. Supports that cannot contain the desired
  equilibrium are discarded.
\item
  \textbf{Construct indifference equations.} For every mixing player,
  expected payoffs from actions in the support must agree. One
  probability per player is recovered from the adding-up restriction, so
  the nonlinear system is expressed in the remaining free probabilities.
\item
  \textbf{Solve the nonlinear system.} In the default method, the free
  probabilities initially lie in interval boxes contained in
  ({[}0,1{]}\^{}d). Interval contraction and subdivision are used to
  search the region systematically.
\item
  \textbf{Newton refinement.} Candidate solution regions are refined to
  point solutions.
\item
  \textbf{Equilibrium verification and consolidation.} Candidate
  probability vectors are checked against equilibrium conditions and
  nearby numerical duplicates are consolidated.
\end{enumerate}

The \texttt{fast} option skips the interval-search stage and instead
starts Newton iteration from randomly drawn points in the relevant
mixed-strategy simplices. This is much faster in many problems, but it
does not provide the same systematic search of the strategy space.

The default procedure is therefore best described as a complete-search
procedure \textbf{intended to locate all equilibria} of the finite game.
As with any numerical procedure, difficult or degenerate games should be
inspected carefully.

    \section*{Digression on Complete Solution of Nonlinear Polynomial
Systems}\label{digression-on-complete-solution-of-nonlinear-polynomial-systems}
\addcontentsline{toc}{section}{Digression on Complete Solution of
Nonlinear Polynomial Systems}

\subsection{Homotopy methods}\label{homotopy-methods}

Homotopy methods solve a nonlinear system by embedding the problem of
interest in a continuous family of related problems. One begins with a
system \(G(x)=0\) whose solutions are known and constructs a homotopy

\[
H(x,t)=(1-t)G(x)+tF(x)=0,
\qquad t\in[0,1],
\]

where \(F(x)=0\) is the system one ultimately wishes to solve. Starting
from each known solution at \(t=0\), a numerical path-following
algorithm traces the corresponding solution as \(t\) moves toward one.
Under appropriate regularity conditions, the endpoints of these paths at
\(t=1\) provide the solutions to the original system.

This approach is particularly attractive for polynomial systems because
algebraic results can often tell us how many solution paths must be
followed. In games, the equilibrium conditions associated with a
particular support form a polynomial system, so homotopy continuation
can be used to trace candidate equilibria systematically. A major
advantage over ordinary Newton iteration is that one is not simply
choosing arbitrary starting points and hoping that they lie in the
basins of attraction of all the relevant roots. The continuation
construction provides a principled collection of paths that can, under
suitable conditions, recover all isolated solutions of the polynomial
system.

The complications are largely practical. Paths can pass close to
singularities, approach one another, or wander through the complex
domain even when the economically relevant solutions are real. Careful
predictor-corrector methods, endgames, and numerical precision are
therefore important. Moreover, polynomial systems may possess many
complex solutions that must be tracked even though only a small subset
correspond to real probability vectors satisfying equilibrium
inequalities. Homotopy methods are consequently powerful and elegant,
but they may devote substantial computation to solutions that are
ultimately irrelevant for a finite-game problem.

\subsection{Algebraic methods}\label{algebraic-methods}

Algebraic methods treat the equilibrium conditions as a system of
polynomial equations and apply tools from computational algebraic
geometry. The best-known examples use Gröbner bases, resultants, or
related elimination techniques. The basic idea is to transform a
multivariate polynomial system into an equivalent system with a more
tractable algebraic structure. In favorable cases, repeated elimination
reduces the problem to one univariate polynomial; its roots can then be
recovered and substituted backward to obtain the remaining variables.

The attraction of the algebraic approach is that it works with the
global structure of the system rather than searching locally from a
collection of initial values. It can reveal how many algebraic solutions
exist, characterize multiplicities, and sometimes deliver exact or
symbolic representations of solutions. This makes it particularly
appealing from a theoretical perspective. Since mixed-strategy
equilibrium conditions in finite games are polynomial, they fit
naturally into this framework, and algebraic methods can in principle
enumerate the solutions associated with a candidate support. The
difficulty is that symbolic calculations can grow extraordinarily
rapidly as the number and degree of the equations increase.
Gröbner-basis computation in particular may suffer from severe
intermediate-expression growth, making problems that look modest in
economic notation computationally demanding. Exact algebraic computation
also produces all algebraic solutions, including complex roots and roots
lying outside the probability simplex, which must subsequently be
filtered. Thus algebraic methods offer considerable theoretical power,
but their computational burden can become substantial as the numbers of
players, actions, and candidate supports increase.

\subsection{Interval methods}\label{interval-methods}

Interval methods take a different approach. Instead of representing an
unknown probability by a floating-point number \(p\), one represents it
initially by an interval known to contain it, for example \(p\in[0,1]\).
Arithmetic operations are then replaced by interval operations
constructed so that the resulting interval is guaranteed to contain
every value that the corresponding real-valued expression could attain.
Applied to a system of equilibrium equations, this makes it possible to
evaluate an entire box of candidate probabilities at once.

The key computational advantage is exclusion. Suppose a system contains
an equation \(f(x)=0\), and interval evaluation over a box \(X\)
produces

\[
0\notin [f](X).
\]

Then no solution can exist anywhere in that box, and the entire region
can be discarded. Boxes that cannot be excluded can be contracted using
interval versions of Newton or related consistency operators; if a box
remains too large, it can be subdivided and the procedure repeated. The
resulting branch-contract-subdivide algorithm progressively eliminates
regions that cannot contain roots while isolating regions that may
contain them. This makes interval methods particularly attractive when
the purpose is not merely to find a root, but to search systematically
for all real roots in a bounded domain. The price of that reliability is
computational cost. Interval evaluations can be conservative because
repeated occurrences of the same variable are treated as though they
varied independently, producing bounds wider than the true range of a
function. This ``dependency problem'' can force additional subdivision
and greatly increase the number of boxes that must be examined.
Nevertheless, the domain of mixed strategies is naturally bounded---the
product of probability simplices---and only real solutions are
economically relevant. Those features make interval methods unusually
well suited to the objective of dagamesolve: search candidate supports
over a known bounded region, rule out regions that cannot contain an
equilibrium, isolate the remaining solutions, and then use an ordinary
point method such as Newton iteration to refine them.

    \section{A short introduction to interval
methods}\label{a-short-introduction-to-interval-methods}

The interval component is what distinguishes the default mixed-strategy
search from simply running Newton's method from a collection of guesses.

Let

\[
[x]=[\underline{x},\overline{x}]
\]

denote an interval known to contain a scalar (x). Basic interval
operations are defined so that the resulting interval contains every
possible real result.

Addition:

\[
[a,b]+[c,d]=[a+c,b+d].
\]

Subtraction:

\[
[a,b]-[c,d]=[a-d,b-c].
\]

Multiplication:

\[
[a,b][c,d]
=
[\min(ac,ad,bc,bd),\max(ac,ad,bc,bd)].
\]

Floating-point implementations use \textbf{outward rounding}: lower
bounds are rounded downward and upper bounds upward so that numerical
rounding does not incorrectly exclude feasible values.

    \subsection{Inclusion functions and
pessimism}\label{inclusion-functions-and-pessimism}

Replacing ordinary arithmetic in an expression by interval arithmetic
produces an inclusion function \([f]([x])\) satisfying

\[
f([x])\subseteq [f]([x]).
\]

The inclusion can be wider than the true range because the same variable
may appear repeatedly in the expression. This is often called the
dependency problem or interval pessimism.

For example, if \(x\in[-1,2]\), the true range of \(x^2\) is \([0,4]\).
Evaluating the algebraically equivalent expression \(x\times x\) by
naive interval multiplication gives

\[
[-1,2]\times[-1,2]=[-2,4],
\]

which contains the true range but is wider.

This apparent conservatism is useful for exclusion. If zero is
\textbf{not} contained in the interval evaluation of an equation over a
box, no root can lie anywhere inside that box.

    \subsection{Branch, contract, and
exclude}\label{branch-contract-and-exclude}

Suppose a scalar root is sought on an initial interval \([x]\).

A generic interval search repeats the following operations:

\begin{enumerate}
\def\labelenumi{\arabic{enumi}.}
\tightlist
\item
  Evaluate an interval inclusion of the equation.
\item
  If \(0\notin[f]([x])\), discard the interval.
\item
  If zero remains possible, attempt to \textbf{contract} the interval.
\item
  If the surviving interval is still too wide, bisect it and place the
  two smaller intervals back on the work list.
\item
  Continue until no unresolved boxes remain.
\item
  Refine sufficiently small candidate boxes with a point method such as
  Newton iteration.
\end{enumerate}

In several dimensions, intervals become boxes and the Jacobian becomes
an interval matrix. A Taylor-type inclusion has the schematic form

\[
\mathbf{f}([\mathbf{x}])
\subseteq
\mathbf{f}(\mathbf{m})
+
[\mathbf{J}]([\mathbf {x}])([\mathbf{x}]-\mathbf{m}),
\]

where \(\mathbf{m}\) is typically a point such as the midpoint of the
box.

The interval solver used by \texttt{dagamesolve} performs
equation-by-equation contraction and subdivision through the companion
\texttt{intsolver} package.

    \section{Example 1: a two-player coordination
game}\label{example-1-a-two-player-coordination-game}

Consider the game

{\def\LTcaptype{none} % do not increment counter
\begin{longtable}[]{@{}lrr@{}}
\toprule\noalign{}
& Player 2: 1 & Player 2: 2 \\
\midrule\noalign{}
\endhead
\bottomrule\noalign{}
\endlastfoot
\textbf{Player 1: 1} & (1,2) & (0,0.1) \\
\textbf{Player 1: 2} & (0,0) & (3,1) \\
\end{longtable}
}

The game has two pure-strategy equilibria and one mixed equilibrium.

The code below constructs the action profiles with
\texttt{dagamestrats}, enters the payoffs in matching profile order, and
solves the game.

    \begin{tcolorbox}[breakable, size=fbox, boxrule=1pt, pad at break*=1mm,colback=cellbackground, colframe=cellborder]
\prompt{In}{incolor}{84}{\boxspacing}
\begin{Verbatim}[commandchars=\\\{\}]
\PY{k}{clear}
\PY{k}{set} obs \PY{l+m}{2}

\PY{k}{gen} game = \PY{l+m}{1}
\PY{k}{gen} acts = \PY{l+m}{2}
dagamestrats acts, group(game)\PY{k}{ generate}(s)

\PY{c+c1}{*}\PY{c+c1}{ }\PY{c+c1}{P}\PY{c+c1}{r}\PY{c+c1}{o}\PY{c+c1}{f}\PY{c+c1}{i}\PY{c+c1}{l}\PY{c+c1}{e}\PY{c+c1}{ }\PY{c+c1}{1}\PY{c+c1}{:}\PY{c+c1}{ }\PY{c+c1}{(}\PY{c+c1}{1}\PY{c+c1}{,}\PY{c+c1}{1}\PY{c+c1}{)}\PY{c+c1}{ }\PY{c+c1}{\PYZhy{}}\PY{c+c1}{\PYZgt{}}\PY{c+c1}{ }\PY{c+c1}{p}\PY{c+c1}{a}\PY{c+c1}{y}\PY{c+c1}{o}\PY{c+c1}{f}\PY{c+c1}{f}\PY{c+c1}{s}\PY{c+c1}{ }\PY{c+c1}{(}\PY{c+c1}{1}\PY{c+c1}{,}\PY{c+c1}{2}\PY{c+c1}{)}
\PY{k}{gen} pay1 = \PY{l+m}{1}
\PY{k}{replace} pay1 = \PY{l+m}{2}\PY{k}{ in} \PY{l+m}{2}

\PY{c+c1}{*}\PY{c+c1}{ }\PY{c+c1}{P}\PY{c+c1}{r}\PY{c+c1}{o}\PY{c+c1}{f}\PY{c+c1}{i}\PY{c+c1}{l}\PY{c+c1}{e}\PY{c+c1}{ }\PY{c+c1}{2}\PY{c+c1}{:}\PY{c+c1}{ }\PY{c+c1}{(}\PY{c+c1}{1}\PY{c+c1}{,}\PY{c+c1}{2}\PY{c+c1}{)}\PY{c+c1}{ }\PY{c+c1}{\PYZhy{}}\PY{c+c1}{\PYZgt{}}\PY{c+c1}{ }\PY{c+c1}{p}\PY{c+c1}{a}\PY{c+c1}{y}\PY{c+c1}{o}\PY{c+c1}{f}\PY{c+c1}{f}\PY{c+c1}{s}\PY{c+c1}{ }\PY{c+c1}{(}\PY{c+c1}{0}\PY{c+c1}{,}\PY{c+c1}{.}\PY{c+c1}{1}\PY{c+c1}{)}
\PY{k}{gen} pay2 = \PY{l+m}{0}
\PY{k}{replace} pay2 = .\PY{l+m}{1}\PY{k}{ in} \PY{l+m}{2}

\PY{c+c1}{*}\PY{c+c1}{ }\PY{c+c1}{P}\PY{c+c1}{r}\PY{c+c1}{o}\PY{c+c1}{f}\PY{c+c1}{i}\PY{c+c1}{l}\PY{c+c1}{e}\PY{c+c1}{ }\PY{c+c1}{3}\PY{c+c1}{:}\PY{c+c1}{ }\PY{c+c1}{(}\PY{c+c1}{2}\PY{c+c1}{,}\PY{c+c1}{1}\PY{c+c1}{)}\PY{c+c1}{ }\PY{c+c1}{\PYZhy{}}\PY{c+c1}{\PYZgt{}}\PY{c+c1}{ }\PY{c+c1}{p}\PY{c+c1}{a}\PY{c+c1}{y}\PY{c+c1}{o}\PY{c+c1}{f}\PY{c+c1}{f}\PY{c+c1}{s}\PY{c+c1}{ }\PY{c+c1}{(}\PY{c+c1}{0}\PY{c+c1}{,}\PY{c+c1}{0}\PY{c+c1}{)}
\PY{k}{gen} pay3 = \PY{l+m}{0}

\PY{c+c1}{*}\PY{c+c1}{ }\PY{c+c1}{P}\PY{c+c1}{r}\PY{c+c1}{o}\PY{c+c1}{f}\PY{c+c1}{i}\PY{c+c1}{l}\PY{c+c1}{e}\PY{c+c1}{ }\PY{c+c1}{4}\PY{c+c1}{:}\PY{c+c1}{ }\PY{c+c1}{(}\PY{c+c1}{2}\PY{c+c1}{,}\PY{c+c1}{2}\PY{c+c1}{)}\PY{c+c1}{ }\PY{c+c1}{\PYZhy{}}\PY{c+c1}{\PYZgt{}}\PY{c+c1}{ }\PY{c+c1}{p}\PY{c+c1}{a}\PY{c+c1}{y}\PY{c+c1}{o}\PY{c+c1}{f}\PY{c+c1}{f}\PY{c+c1}{s}\PY{c+c1}{ }\PY{c+c1}{(}\PY{c+c1}{3}\PY{c+c1}{,}\PY{c+c1}{1}\PY{c+c1}{)}
\PY{k}{gen} pay4 = \PY{l+m}{3}
\PY{k}{replace} pay4 = \PY{l+m}{1}\PY{k}{ in} \PY{l+m}{2}

\PY{k}{list} game acts s1\PY{o}{\PYZhy{}}s4 pay1\PY{o}{\PYZhy{}}pay4, noobs

dagamesolve s1\PY{o}{\PYZhy{}}s4, group(game) payoffs(pay1\PY{o}{\PYZhy{}}pay4) equilibria(eq)

\PY{k}{return}\PY{k}{ list}

\PY{k}{format} eq\PYZus{}\PY{o}{*} \PY{n+nx}{\PYZpc{}6.3f}
\PY{k}{format} eq\PYZus{}count \PY{n+nx}{\PYZpc{}3.0f}

\PY{k}{format} eq\PYZus{}\PY{o}{*} \PY{n+nx}{\PYZpc{}6.3f}

\PY{k}{list} eq\PYZus{}1\PY{o}{*} eq\PYZus{}rets\PYZus{}1, noobs sep(\PY{l+m}{0})
\PY{k}{list} eq\PYZus{}2\PY{o}{*} eq\PYZus{}rets\PYZus{}2, noobs sep(\PY{l+m}{0})
\PY{k}{list} eq\PYZus{}3\PY{o}{*} eq\PYZus{}rets\PYZus{}3, noobs sep(\PY{l+m}{0})
\end{Verbatim}
\end{tcolorbox}

    \begin{Verbatim}[commandchars=\\\{\}]
Number of observations (\_N) was 0, now 2.
(1 real change made)
(1 real change made)
(1 real change made)

  +-------------------------------------------------------------+
  | game   acts   s1   s2   s3   s4   pay1   pay2   pay3   pay4 |
  |-------------------------------------------------------------|
  |    1      2    1    1    2    2      1      0      0      3 |
  |    1      2    1    2    1    2      2     .1      0      1 |
  +-------------------------------------------------------------+

scalars:
              r(sdeqs) =  .
             r(mineqs) =  3
             r(maxeqs) =  3
            r(meaneqs) =  3
              r(games) =  1
          r(tolerance) =  1.00000000000e-12
             r(digits) =  1.00000000000e-08
              r(mnits) =  100
             r(coltol) =  .0001
            r(mbwidth) =  .0001
       r(totalactions) =  4

macros:
             r(method) : "complete"

matrices:
      r(playeractions) :  1 x 2

  +------------------------------------------------------+
  | eq\_1\_1\_1   eq\_1\_1\_2   eq\_1\_2\_1   eq\_1\_2\_2   eq\_ret\textasciitilde{}1 |
  |------------------------------------------------------|
  |        .          .          .          .      1.000 |
  |    1.000      0.000      1.000      0.000      2.000 |
  +------------------------------------------------------+

  +------------------------------------------------------+
  | eq\_2\_1\_1   eq\_2\_1\_2   eq\_2\_2\_1   eq\_2\_2\_2   eq\_ret\textasciitilde{}2 |
  |------------------------------------------------------|
  |        .          .          .          .      3.000 |
  |    0.000      1.000      0.000      1.000      1.000 |
  +------------------------------------------------------+

  +------------------------------------------------------+
  | eq\_3\_1\_1   eq\_3\_1\_2   eq\_3\_2\_1   eq\_3\_2\_2   eq\_ret\textasciitilde{}3 |
  |------------------------------------------------------|
  |        .          .          .          .      0.750 |
  |    0.345      0.655      0.750      0.250      0.690 |
  +------------------------------------------------------+
    \end{Verbatim}

    If \texttt{equilibria(eq)} is specified, strategy probabilities are
stored in variables of the form

\[
\texttt{eq\_q\_p\_a},
\]

where \texttt{q} indexes the equilibrium, \texttt{p} the player, and
\texttt{a} the action. The complete strategy profile is stored on the
last player observation for the game. Variables \texttt{eq\_rets\_q}
contain the player-specific equilibrium payoffs, and \texttt{eq\_count}
records the number of equilibria found.

    \section{Example 2: matching pennies}\label{example-2-matching-pennies}

Matching pennies has no pure-strategy equilibrium. Its unique Nash
equilibrium is fully mixed, with each player placing probability (1/2)
on each action.

This is a useful minimal example because a solver restricted to pure
strategies would report no equilibrium, while a mixed-strategy solver
should recover the interior solution.

    \begin{tcolorbox}[breakable, size=fbox, boxrule=1pt, pad at break*=1mm,colback=cellbackground, colframe=cellborder]
\prompt{In}{incolor}{87}{\boxspacing}
\begin{Verbatim}[commandchars=\\\{\}]
\PY{k}{clear}
\PY{k}{set} obs \PY{l+m}{2}

\PY{k}{gen} game = \PY{l+m}{1}
\PY{k}{gen} acts = \PY{l+m}{2}
dagamestrats acts, group(game)\PY{k}{ generate}(s)

\PY{c+c1}{*}\PY{c+c1}{ }\PY{c+c1}{P}\PY{c+c1}{l}\PY{c+c1}{a}\PY{c+c1}{y}\PY{c+c1}{e}\PY{c+c1}{r}\PY{c+c1}{ }\PY{c+c1}{1}\PY{c+c1}{:}\PY{c+c1}{ }\PY{c+c1}{ }\PY{c+c1}{1}\PY{c+c1}{,}\PY{c+c1}{ }\PY{c+c1}{\PYZhy{}}\PY{c+c1}{1}\PY{c+c1}{,}\PY{c+c1}{ }\PY{c+c1}{\PYZhy{}}\PY{c+c1}{1}\PY{c+c1}{,}\PY{c+c1}{ }\PY{c+c1}{ }\PY{c+c1}{1}
\PY{c+c1}{*}\PY{c+c1}{ }\PY{c+c1}{P}\PY{c+c1}{l}\PY{c+c1}{a}\PY{c+c1}{y}\PY{c+c1}{e}\PY{c+c1}{r}\PY{c+c1}{ }\PY{c+c1}{2}\PY{c+c1}{:}\PY{c+c1}{ }\PY{c+c1}{\PYZhy{}}\PY{c+c1}{1}\PY{c+c1}{,}\PY{c+c1}{ }\PY{c+c1}{ }\PY{c+c1}{1}\PY{c+c1}{,}\PY{c+c1}{ }\PY{c+c1}{ }\PY{c+c1}{1}\PY{c+c1}{,}\PY{c+c1}{ }\PY{c+c1}{\PYZhy{}}\PY{c+c1}{1}
\PY{k}{gen} pay1 = \PY{n+nf}{cond}(\PYZus{}n\PY{o}{==}\PY{l+m}{1},  \PY{l+m}{1}, \PY{o}{\PYZhy{}}\PY{l+m}{1})
\PY{k}{gen} pay2 = \PY{n+nf}{cond}(\PYZus{}n\PY{o}{==}\PY{l+m}{1}, \PY{o}{\PYZhy{}}\PY{l+m}{1},  \PY{l+m}{1})
\PY{k}{gen} pay3 = \PY{n+nf}{cond}(\PYZus{}n\PY{o}{==}\PY{l+m}{1}, \PY{o}{\PYZhy{}}\PY{l+m}{1},  \PY{l+m}{1})
\PY{k}{gen} pay4 = \PY{n+nf}{cond}(\PYZus{}n\PY{o}{==}\PY{l+m}{1},  \PY{l+m}{1}, \PY{o}{\PYZhy{}}\PY{l+m}{1})

dagamesolve s1\PY{o}{\PYZhy{}}s4, group(game) payoffs(pay1\PY{o}{\PYZhy{}}pay4) equilibria(eq)

\PY{k}{list}\PY{k}{ eq}\PY{o}{*}, noobs
\PY{k}{return}\PY{k}{ list}
\end{Verbatim}
\end{tcolorbox}

    \begin{Verbatim}[commandchars=\\\{\}]
Number of observations (\_N) was 0, now 2.

  +-----------------------------------------------------------------+
  | eq\_ret\textasciitilde{}1   eq\_1\_1\_1   eq\_1\_1\_2   eq\_1\_2\_1   eq\_1\_2\_2   eq\_count |
  |-----------------------------------------------------------------|
  |        0          .          .          .          .          . |
  |        0         .5         .5         .5         .5          1 |
  +-----------------------------------------------------------------+

scalars:
              r(sdeqs) =  .
             r(mineqs) =  1
             r(maxeqs) =  1
            r(meaneqs) =  1
              r(games) =  1
          r(tolerance) =  1.00000000000e-12
             r(digits) =  1.00000000000e-08
              r(mnits) =  100
             r(coltol) =  .0001
            r(mbwidth) =  .0001
       r(totalactions) =  4

macros:
             r(method) : "complete"

matrices:
      r(playeractions) :  1 x 2
    \end{Verbatim}

    \section{Example 3: a three-player coordination
game}\label{example-3-a-three-player-coordination-game}

The earlier presentation used the following three-player game to
illustrate the list representation. In the presentation actions were
numbered 0 and 1; here they are numbered 1 and 2 to match
\texttt{dagamestrats}.

Only the two unanimous profiles have nonzero payoffs:

{[} \Pi(1,1,1)=(2,1,4), \qquad \Pi(2,2,2)=(1.2,2,3), {]}

and the six remaining profiles give all players zero.

The presentation identifies three equilibria in this game. The following
cell constructs and solves it directly.

    \begin{tcolorbox}[breakable, size=fbox, boxrule=1pt, pad at break*=1mm,colback=cellbackground, colframe=cellborder]
\prompt{In}{incolor}{89}{\boxspacing}
\begin{Verbatim}[commandchars=\\\{\}]
\PY{k}{clear}
\PY{k}{set} obs \PY{l+m}{3}

\PY{k}{gen} game = \PY{l+m}{1}
\PY{k}{gen} acts = \PY{l+m}{2}
dagamestrats acts, group(game)\PY{k}{ generate}(s)

\PY{k}{forvalues} j = \PY{l+m}{1}\PY{o}{/}\PY{l+m}{8} \PYZob{}
\PY{k}{    gen} pay\PY{n+nv}{`}\PY{n+nv}{j}\PY{n+nv}{\PYZsq{}} = \PY{l+m}{0}
\PYZcb{}

\PY{c+c1}{*}\PY{c+c1}{ }\PY{c+c1}{U}\PY{c+c1}{n}\PY{c+c1}{a}\PY{c+c1}{n}\PY{c+c1}{i}\PY{c+c1}{m}\PY{c+c1}{o}\PY{c+c1}{u}\PY{c+c1}{s}\PY{c+c1}{ }\PY{c+c1}{a}\PY{c+c1}{c}\PY{c+c1}{t}\PY{c+c1}{i}\PY{c+c1}{o}\PY{c+c1}{n}\PY{c+c1}{ }\PY{c+c1}{1}
\PY{k}{replace} pay1 = \PY{l+m}{2}\PY{k}{ in} \PY{l+m}{1}
\PY{k}{replace} pay1 = \PY{l+m}{1}\PY{k}{ in} \PY{l+m}{2}
\PY{k}{replace} pay1 = \PY{l+m}{4}\PY{k}{ in} \PY{l+m}{3}

\PY{c+c1}{*}\PY{c+c1}{ }\PY{c+c1}{U}\PY{c+c1}{n}\PY{c+c1}{a}\PY{c+c1}{n}\PY{c+c1}{i}\PY{c+c1}{m}\PY{c+c1}{o}\PY{c+c1}{u}\PY{c+c1}{s}\PY{c+c1}{ }\PY{c+c1}{a}\PY{c+c1}{c}\PY{c+c1}{t}\PY{c+c1}{i}\PY{c+c1}{o}\PY{c+c1}{n}\PY{c+c1}{ }\PY{c+c1}{2}
\PY{k}{replace} pay8 = \PY{l+m}{1.2}\PY{k}{ in} \PY{l+m}{1}
\PY{k}{replace} pay8 = \PY{l+m}{2}  \PY{k}{ in} \PY{l+m}{2}
\PY{k}{replace} pay8 = \PY{l+m}{3}  \PY{k}{ in} \PY{l+m}{3}

\PY{k}{list} s1\PY{o}{\PYZhy{}}s8, noobs
\PY{k}{list} pay1\PY{o}{\PYZhy{}}pay8, noobs

dagamesolve s1\PY{o}{\PYZhy{}}s8, group(game) payoffs(pay1\PY{o}{\PYZhy{}}pay8) equilibria(eq)

\PY{k}{return}\PY{k}{ list}

\PY{k}{list} eq\PYZus{}1\PY{o}{*} eq\PYZus{}rets\PYZus{}1, noobs sep(\PY{l+m}{0})
\PY{k}{list} eq\PYZus{}2\PY{o}{*} eq\PYZus{}rets\PYZus{}2, noobs sep(\PY{l+m}{0})
\PY{k}{list} eq\PYZus{}3\PY{o}{*} eq\PYZus{}rets\PYZus{}3, noobs sep(\PY{l+m}{0})
\end{Verbatim}
\end{tcolorbox}

    \begin{Verbatim}[commandchars=\\\{\}]
Number of observations (\_N) was 0, now 3.
(1 real change made)
(1 real change made)
(1 real change made)
(1 real change made)
(1 real change made)
(1 real change made)

  +---------------------------------------+
  | s1   s2   s3   s4   s5   s6   s7   s8 |
  |---------------------------------------|
  |  1    1    1    1    2    2    2    2 |
  |  1    1    2    2    1    1    2    2 |
  |  1    2    1    2    1    2    1    2 |
  +---------------------------------------+

  +-------------------------------------------------------+
  | pay1   pay2   pay3   pay4   pay5   pay6   pay7   pay8 |
  |-------------------------------------------------------|
  |    2      0      0      0      0      0      0    1.2 |
  |    1      0      0      0      0      0      0      2 |
  |    4      0      0      0      0      0      0      3 |
  +-------------------------------------------------------+

scalars:
              r(sdeqs) =  .
             r(mineqs) =  3
             r(maxeqs) =  3
            r(meaneqs) =  3
              r(games) =  1
          r(tolerance) =  1.00000000000e-12
             r(digits) =  1.00000000000e-08
              r(mnits) =  100
             r(coltol) =  .0001
            r(mbwidth) =  .0001
       r(totalactions) =  6

macros:
             r(method) : "complete"

matrices:
      r(playeractions) :  1 x 3

  +----------------------------------------------------------------------------+
  | eq\_1\_1\_1   eq\_1\_1\_2   eq\_1\_2\_1   eq\_1\_2\_2   eq\_1\_3\_1   eq\_1\_3\_2   eq\_ret\textasciitilde{}1 |
  |----------------------------------------------------------------------------|
  |        .          .          .          .          .          .          2 |
  |        .          .          .          .          .          .          1 |
  |        1          0          1          0          1          0          4 |
  +----------------------------------------------------------------------------+

  +----------------------------------------------------------------------------+
  | eq\_2\_1\_1   eq\_2\_1\_2   eq\_2\_2\_1   eq\_2\_2\_2   eq\_2\_3\_1   eq\_2\_3\_2   eq\_ret\textasciitilde{}2 |
  |----------------------------------------------------------------------------|
  |        .          .          .          .          .          .        1.2 |
  |        .          .          .          .          .          .          2 |
  |        0          1          0          1          0          1          3 |
  +----------------------------------------------------------------------------+

  +----------------------------------------------------------------------------+
  | eq\_3\_1\_1   eq\_3\_1\_2   eq\_3\_2\_1   eq\_3\_2\_2   eq\_3\_3\_1   eq\_3\_3\_2   eq\_ret\textasciitilde{}3 |
  |----------------------------------------------------------------------------|
  |        .          .          .          .          .          .   .3593619 |
  |        .          .          .          .          .          .   .3421113 |
  | .6125741   .3874259   .3217312   .6782688   .5584816   .4415184   .7883368 |
  +----------------------------------------------------------------------------+
    \end{Verbatim}

    \section{The distinctive feature: solving a sequence of
games}\label{the-distinctive-feature-solving-a-sequence-of-games}

For empirical work, it is often more useful to solve a family of related
games than a single game.

Consider a two-firm entry game. Each player chooses \textbf{Enter} or
\textbf{Stay Out}. Normalize the payoff from staying out to zero. If a
firm enters alone, its payoff is (\pi\textgreater0). If both firms
enter, each receives

\[
\pi-\delta.
\]

Thus:

\begin{itemize}
\tightlist
\item
  when \(\delta<\pi\), entry is a dominant action and the unique
  equilibrium is joint entry;
\item
  when \(\delta>\pi\), the game has two asymmetric pure equilibria and a
  mixed equilibrium;
\item
  the knife-edge case \(\delta=\pi\) is deliberately omitted below.
\end{itemize}

Instead of solving ten games one at a time, we place all ten in one
Stata dataset and call \texttt{dagamesolve} once.

    The equilibrium count should change as \(\delta\) crosses \(\pi\).
Because the equilibrium count is stored on the last observation of each
game, a simple plot can summarize the sequence.

    \begin{tcolorbox}[breakable, size=fbox, boxrule=1pt, pad at break*=1mm,colback=cellbackground, colframe=cellborder]
\prompt{In}{incolor}{92}{\boxspacing}
\begin{Verbatim}[commandchars=\\\{\}]
\PY{k}{clear}
\PY{k}{set} obs \PY{l+m}{20}

\PY{k}{gen} game = \PY{n+nf}{ceil}(\PYZus{}n\PY{o}{/}\PY{l+m}{2})
\PY{k}{bysort} game:\PY{k}{ gen} player = \PYZus{}n

\PY{k}{gen} acts = \PY{l+m}{2}
dagamestrats acts, group(game)\PY{k}{ generate}(s)

\PY{c+c1}{*}\PY{c+c1}{ }\PY{c+c1}{C}\PY{c+c1}{o}\PY{c+c1}{m}\PY{c+c1}{m}\PY{c+c1}{o}\PY{c+c1}{n}\PY{c+c1}{ }\PY{c+c1}{m}\PY{c+c1}{o}\PY{c+c1}{n}\PY{c+c1}{o}\PY{c+c1}{p}\PY{c+c1}{o}\PY{c+c1}{l}\PY{c+c1}{y}\PY{c+c1}{\PYZhy{}}\PY{c+c1}{e}\PY{c+c1}{n}\PY{c+c1}{t}\PY{c+c1}{r}\PY{c+c1}{y}\PY{c+c1}{ }\PY{c+c1}{p}\PY{c+c1}{a}\PY{c+c1}{y}\PY{c+c1}{o}\PY{c+c1}{f}\PY{c+c1}{f}
\PY{k}{gen} pi = \PY{l+m}{1}

\PY{c+c1}{*}\PY{c+c1}{ }\PY{c+c1}{d}\PY{c+c1}{e}\PY{c+c1}{l}\PY{c+c1}{t}\PY{c+c1}{a}\PY{c+c1}{ }\PY{c+c1}{=}\PY{c+c1}{ }\PY{c+c1}{.}\PY{c+c1}{5}\PY{c+c1}{,}\PY{c+c1}{.}\PY{c+c1}{6}\PY{c+c1}{,}\PY{c+c1}{.}\PY{c+c1}{7}\PY{c+c1}{,}\PY{c+c1}{.}\PY{c+c1}{8}\PY{c+c1}{,}\PY{c+c1}{.}\PY{c+c1}{9}\PY{c+c1}{,}\PY{c+c1}{1}\PY{c+c1}{.}\PY{c+c1}{1}\PY{c+c1}{,}\PY{c+c1}{1}\PY{c+c1}{.}\PY{c+c1}{2}\PY{c+c1}{,}\PY{c+c1}{1}\PY{c+c1}{.}\PY{c+c1}{3}\PY{c+c1}{,}\PY{c+c1}{1}\PY{c+c1}{.}\PY{c+c1}{4}\PY{c+c1}{,}\PY{c+c1}{1}\PY{c+c1}{.}\PY{c+c1}{5}
\PY{k}{gen} delta = .\PY{l+m}{4} \PY{o}{+} .\PY{l+m}{1}\PY{o}{*}game
\PY{k}{replace} delta = delta \PY{o}{+} .\PY{l+m}{1}\PY{k}{ if} delta \PY{o}{\PYZgt{}=} \PY{l+m}{1}

\PY{c+c1}{*}\PY{c+c1}{ }\PY{c+c1}{s}\PY{c+c1}{1}\PY{c+c1}{ }\PY{c+c1}{=}\PY{c+c1}{ }\PY{c+c1}{(}\PY{c+c1}{E}\PY{c+c1}{n}\PY{c+c1}{t}\PY{c+c1}{e}\PY{c+c1}{r}\PY{c+c1}{,}\PY{c+c1}{E}\PY{c+c1}{n}\PY{c+c1}{t}\PY{c+c1}{e}\PY{c+c1}{r}\PY{c+c1}{)}
\PY{k}{gen} pay1 = pi \PY{o}{\PYZhy{}} delta

\PY{c+c1}{*}\PY{c+c1}{ }\PY{c+c1}{s}\PY{c+c1}{2}\PY{c+c1}{ }\PY{c+c1}{=}\PY{c+c1}{ }\PY{c+c1}{(}\PY{c+c1}{E}\PY{c+c1}{n}\PY{c+c1}{t}\PY{c+c1}{e}\PY{c+c1}{r}\PY{c+c1}{,}\PY{c+c1}{S}\PY{c+c1}{t}\PY{c+c1}{a}\PY{c+c1}{y}\PY{c+c1}{O}\PY{c+c1}{u}\PY{c+c1}{t}\PY{c+c1}{)}
\PY{k}{gen} pay2 = \PY{n+nf}{cond}(player\PY{o}{==}\PY{l+m}{1}, pi, \PY{l+m}{0})

\PY{c+c1}{*}\PY{c+c1}{ }\PY{c+c1}{s}\PY{c+c1}{3}\PY{c+c1}{ }\PY{c+c1}{=}\PY{c+c1}{ }\PY{c+c1}{(}\PY{c+c1}{S}\PY{c+c1}{t}\PY{c+c1}{a}\PY{c+c1}{y}\PY{c+c1}{O}\PY{c+c1}{u}\PY{c+c1}{t}\PY{c+c1}{,}\PY{c+c1}{E}\PY{c+c1}{n}\PY{c+c1}{t}\PY{c+c1}{e}\PY{c+c1}{r}\PY{c+c1}{)}
\PY{k}{gen} pay3 = \PY{n+nf}{cond}(player\PY{o}{==}\PY{l+m}{1}, \PY{l+m}{0}, pi)

\PY{c+c1}{*}\PY{c+c1}{ }\PY{c+c1}{s}\PY{c+c1}{4}\PY{c+c1}{ }\PY{c+c1}{=}\PY{c+c1}{ }\PY{c+c1}{(}\PY{c+c1}{S}\PY{c+c1}{t}\PY{c+c1}{a}\PY{c+c1}{y}\PY{c+c1}{O}\PY{c+c1}{u}\PY{c+c1}{t}\PY{c+c1}{,}\PY{c+c1}{S}\PY{c+c1}{t}\PY{c+c1}{a}\PY{c+c1}{y}\PY{c+c1}{O}\PY{c+c1}{u}\PY{c+c1}{t}\PY{c+c1}{)}
\PY{k}{gen} pay4 = \PY{l+m}{0}

dagamesolve s1\PY{o}{\PYZhy{}}s4, group(game) payoffs(pay1\PY{o}{\PYZhy{}}pay4) equilibria(eq)

\PY{k}{list} game delta eq\PYZus{}count\PY{k}{ if} player\PY{o}{==}\PY{l+m}{2}, noobs sep(\PY{l+m}{0})
\end{Verbatim}
\end{tcolorbox}

    \begin{Verbatim}[commandchars=\\\{\}]
Number of observations (\_N) was 0, now 20.
(10 real changes made)

  +-------------------------+
  | game   delta   eq\_count |
  |-------------------------|
  |    1      .5          1 |
  |    2      .6          1 |
  |    3      .7          1 |
  |    4      .8          1 |
  |    5      .9          1 |
  |    6     1.1          3 |
  |    7     1.2          3 |
  |    8     1.3          3 |
  |    9     1.4          3 |
  |   10     1.5          3 |
  +-------------------------+
    \end{Verbatim}

    \begin{tcolorbox}[breakable, size=fbox, boxrule=1pt, pad at break*=1mm,colback=cellbackground, colframe=cellborder]
\prompt{In}{incolor}{93}{\boxspacing}
\begin{Verbatim}[commandchars=\\\{\}]
\PY{k}{sort} delta

\PY{k}{twoway} connected eq\PYZus{}count delta\PY{k}{ if} player\PY{o}{==}\PY{l+m}{2}, \PY{c+cs}{///}
    ytitle(\PY{l+s}{\PYZdq{}}\PY{l+s}{N}\PY{l+s}{u}\PY{l+s}{m}\PY{l+s}{b}\PY{l+s}{e}\PY{l+s}{r}\PY{l+s}{ }\PY{l+s}{o}\PY{l+s}{f}\PY{l+s}{ }\PY{l+s}{N}\PY{l+s}{a}\PY{l+s}{s}\PY{l+s}{h}\PY{l+s}{ }\PY{l+s}{e}\PY{l+s}{q}\PY{l+s}{u}\PY{l+s}{i}\PY{l+s}{l}\PY{l+s}{i}\PY{l+s}{b}\PY{l+s}{r}\PY{l+s}{i}\PY{l+s}{a}\PY{l+s}{\PYZdq{}}) \PY{c+cs}{///}
    xtitle(\PY{l+s}{\PYZdq{}}\PY{l+s}{C}\PY{l+s}{o}\PY{l+s}{m}\PY{l+s}{p}\PY{l+s}{e}\PY{l+s}{t}\PY{l+s}{i}\PY{l+s}{t}\PY{l+s}{i}\PY{l+s}{o}\PY{l+s}{n}\PY{l+s}{ }\PY{l+s}{p}\PY{l+s}{a}\PY{l+s}{r}\PY{l+s}{a}\PY{l+s}{m}\PY{l+s}{e}\PY{l+s}{t}\PY{l+s}{e}\PY{l+s}{r}\PY{l+s}{ }\PY{l+s}{d}\PY{l+s}{e}\PY{l+s}{l}\PY{l+s}{t}\PY{l+s}{a}\PY{l+s}{\PYZdq{}}) \PY{c+cs}{///}
    title(\PY{l+s}{\PYZdq{}}\PY{l+s}{E}\PY{l+s}{q}\PY{l+s}{u}\PY{l+s}{i}\PY{l+s}{l}\PY{l+s}{i}\PY{l+s}{b}\PY{l+s}{r}\PY{l+s}{i}\PY{l+s}{u}\PY{l+s}{m}\PY{l+s}{ }\PY{l+s}{m}\PY{l+s}{u}\PY{l+s}{l}\PY{l+s}{t}\PY{l+s}{i}\PY{l+s}{p}\PY{l+s}{l}\PY{l+s}{i}\PY{l+s}{c}\PY{l+s}{i}\PY{l+s}{t}\PY{l+s}{y}\PY{l+s}{ }\PY{l+s}{i}\PY{l+s}{n}\PY{l+s}{ }\PY{l+s}{a}\PY{l+s}{ }\PY{l+s}{s}\PY{l+s}{e}\PY{l+s}{q}\PY{l+s}{u}\PY{l+s}{e}\PY{l+s}{n}\PY{l+s}{c}\PY{l+s}{e}\PY{l+s}{ }\PY{l+s}{o}\PY{l+s}{f}\PY{l+s}{ }\PY{l+s}{e}\PY{l+s}{n}\PY{l+s}{t}\PY{l+s}{r}\PY{l+s}{y}\PY{l+s}{ }\PY{l+s}{g}\PY{l+s}{a}\PY{l+s}{m}\PY{l+s}{e}\PY{l+s}{s}\PY{l+s}{\PYZdq{}})
\end{Verbatim}
\end{tcolorbox}

    \begin{center}
    \adjustimage{max size={0.9\linewidth}{0.9\paperheight}}{dagamesolve_paper_files/dagamesolve_paper_29_0.png}
    \end{center}
    { \hspace*{\fill} \\}
    
    This example captures an important intended use of the software. A
researcher can generate game-specific payoffs from covariates or
parameter values, solve the entire sequence, and then feed equilibrium
objects into subsequent analysis without moving between Stata and a
stand-alone game-theory program.

    \section{\texorpdfstring{Complete search, \texttt{fast}, and
\texttt{pureonly}}{Complete search, fast, and pureonly}}\label{complete-search-fast-and-pureonly}

The default mixed-strategy method uses interval search followed by
Newton refinement. Two useful alternatives are available.

\subsection{\texorpdfstring{\texttt{fast}}{fast}}\label{fast}

With \texttt{fast}, the interval-search stage is skipped. For each
support, random points are drawn from the relevant mixed-strategy
simplices and used as starting values for Newton iteration.
\texttt{randompoints(\#)} controls the number of starts.

This may be dramatically faster, but it can miss equilibria if none of
the random starting points fall in a useful basin of attraction.

\subsection{\texorpdfstring{\texttt{pureonly}}{pureonly}}\label{pureonly}

With \texttt{pureonly}, no mixed-strategy systems are solved. This is
useful when only pure equilibria are substantively relevant or when one
wants a quick preliminary diagnostic.

    \begin{tcolorbox}[breakable, size=fbox, boxrule=1pt, pad at break*=1mm,colback=cellbackground, colframe=cellborder]
\prompt{In}{incolor}{96}{\boxspacing}
\begin{Verbatim}[commandchars=\\\{\}]
\PY{k}{set} seed \PY{l+m}{5150}

dagamesolve s1\PY{o}{\PYZhy{}}s4, group(game) payoffs(pay1\PY{o}{\PYZhy{}}pay4) \PY{c+cs}{///}
    equilibria(ef) fast randompoints(\PY{l+m}{200})

dagamesolve s1\PY{o}{\PYZhy{}}s4, group(game) payoffs(pay1\PY{o}{\PYZhy{}}pay4) \PY{c+cs}{///}
    equilibria(ep) pureonly

\PY{k}{list} game delta eq\PYZus{}count ef\PYZus{}count ep\PYZus{}count\PY{k}{ if} player\PY{o}{==}\PY{l+m}{2}, noobs sep(\PY{l+m}{0})
\end{Verbatim}
\end{tcolorbox}

    \begin{Verbatim}[commandchars=\\\{\}]

  +-----------------------------------------------+
  | game   delta   eq\_count   ef\_count   ep\_count |
  |-----------------------------------------------|
  |    1      .5          1          1          1 |
  |    2      .6          1          1          1 |
  |    3      .7          1          1          1 |
  |    4      .8          1          1          1 |
  |    5      .9          1          1          1 |
  |    6     1.1          3          3          2 |
  |    7     1.2          3          3          2 |
  |    8     1.3          3          3          2 |
  |    9     1.4          3          3          2 |
  |   10     1.5          3          3          2 |
  +-----------------------------------------------+
    \end{Verbatim}

    For this simple sequence one would expect the default and sufficiently
thorough \texttt{fast} searches to agree. The \texttt{pureonly} count
differs when a mixed equilibrium is present. Agreement in a simple
example should not be interpreted as a general completeness guarantee
for \texttt{fast}.

    \section{Numerical controls}\label{numerical-controls}

The default complete-search settings are:

{\def\LTcaptype{none} % do not increment counter
\begin{longtable}[]{@{}
  >{\raggedright\arraybackslash}p{(\linewidth - 4\tabcolsep) * \real{0.3000}}
  >{\raggedleft\arraybackslash}p{(\linewidth - 4\tabcolsep) * \real{0.4000}}
  >{\raggedright\arraybackslash}p{(\linewidth - 4\tabcolsep) * \real{0.3000}}@{}}
\toprule\noalign{}
\begin{minipage}[b]{\linewidth}\raggedright
Option
\end{minipage} & \begin{minipage}[b]{\linewidth}\raggedleft
Default
\end{minipage} & \begin{minipage}[b]{\linewidth}\raggedright
Role
\end{minipage} \\
\midrule\noalign{}
\endhead
\bottomrule\noalign{}
\endlastfoot
\texttt{mbwidth(\#)} & \texttt{1e-4} & Target maximum interval-box width
before refinement \\
\texttt{tolerance(\#)} & \texttt{1e-12} & Newton convergence
tolerance \\
\texttt{coltol(\#)} & \texttt{1e-4} & Tolerance for consolidating
numerically nearby solutions \\
\texttt{digits(\#)} & \texttt{1e-8} & Precision parameter used in
interval outward rounding \\
\texttt{mnits(\#)} & \texttt{100} & Maximum Newton iterations \\
\end{longtable}
}

For the fast search,

{\def\LTcaptype{none} % do not increment counter
\begin{longtable}[]{@{}
  >{\raggedright\arraybackslash}p{(\linewidth - 4\tabcolsep) * \real{0.3000}}
  >{\raggedleft\arraybackslash}p{(\linewidth - 4\tabcolsep) * \real{0.4000}}
  >{\raggedright\arraybackslash}p{(\linewidth - 4\tabcolsep) * \real{0.3000}}@{}}
\toprule\noalign{}
\begin{minipage}[b]{\linewidth}\raggedright
Option
\end{minipage} & \begin{minipage}[b]{\linewidth}\raggedleft
Default
\end{minipage} & \begin{minipage}[b]{\linewidth}\raggedright
Role
\end{minipage} \\
\midrule\noalign{}
\endhead
\bottomrule\noalign{}
\endlastfoot
\texttt{randompoints(\#)} & \texttt{50} & Random Newton starting points
for each candidate support \\
\texttt{tolerance(\#)} & \texttt{1e-12} & Newton convergence
tolerance \\
\texttt{mnits(\#)} & \texttt{100} & Maximum Newton iterations \\
\end{longtable}
}

These parameters control different parts of the numerical procedure and
should not be interpreted as interchangeable measures of ``accuracy.''
In particular, \texttt{mbwidth()} concerns the interval search,
\texttt{tolerance()} concerns Newton iteration, \texttt{coltol()}
concerns consolidation of candidate solutions, and \texttt{digits()}
concerns interval arithmetic.

    \section{A sequence-of-games regression
test}\label{a-sequence-of-games-regression-test}

The package's test suite contains a useful compact example: three
two-player games are stacked in one dataset.

\begin{enumerate}
\def\labelenumi{\arabic{enumi}.}
\tightlist
\item
  a coordination game with three equilibria;
\item
  a dominant-strategy game with one equilibrium;
\item
  matching pennies with one mixed equilibrium.
\end{enumerate}

The example below is adapted directly from the current regression test
and is useful both pedagogically and as a reproducibility check.

    \begin{tcolorbox}[breakable, size=fbox, boxrule=1pt, pad at break*=1mm,colback=cellbackground, colframe=cellborder]
\prompt{In}{incolor}{100}{\boxspacing}
\begin{Verbatim}[commandchars=\\\{\}]
\PY{k}{clear}
\PY{k}{set} obs \PY{l+m}{6}

\PY{k}{gen} game = \PY{n+nf}{ceil}(\PYZus{}n\PY{o}{/}\PY{l+m}{2})
\PY{k}{bysort} game:\PY{k}{ gen} player = \PYZus{}n

\PY{k}{gen} acts = \PY{l+m}{2}
dagamestrats acts, group(game)\PY{k}{ generate}(s)

\PY{k}{forvalues} j = \PY{l+m}{1}\PY{o}{/}\PY{l+m}{4} \PYZob{}
\PY{k}{    gen} pay\PY{n+nv}{`}\PY{n+nv}{j}\PY{n+nv}{\PYZsq{}} = .
\PYZcb{}

\PY{c+c1}{*}\PY{c+c1}{ }\PY{c+c1}{G}\PY{c+c1}{a}\PY{c+c1}{m}\PY{c+c1}{e}\PY{c+c1}{ }\PY{c+c1}{1}\PY{c+c1}{:}\PY{c+c1}{ }\PY{c+c1}{c}\PY{c+c1}{o}\PY{c+c1}{o}\PY{c+c1}{r}\PY{c+c1}{d}\PY{c+c1}{i}\PY{c+c1}{n}\PY{c+c1}{a}\PY{c+c1}{t}\PY{c+c1}{i}\PY{c+c1}{o}\PY{c+c1}{n}
\PY{k}{replace} pay1 = \PY{l+m}{1}\PY{k}{ if} game\PY{o}{==}\PY{l+m}{1}
\PY{k}{replace} pay2 = \PY{l+m}{0}\PY{k}{ if} game\PY{o}{==}\PY{l+m}{1}
\PY{k}{replace} pay3 = \PY{l+m}{0}\PY{k}{ if} game\PY{o}{==}\PY{l+m}{1}
\PY{k}{replace} pay4 = \PY{l+m}{1}\PY{k}{ if} game\PY{o}{==}\PY{l+m}{1}

\PY{c+c1}{*}\PY{c+c1}{ }\PY{c+c1}{G}\PY{c+c1}{a}\PY{c+c1}{m}\PY{c+c1}{e}\PY{c+c1}{ }\PY{c+c1}{2}\PY{c+c1}{:}\PY{c+c1}{ }\PY{c+c1}{a}\PY{c+c1}{c}\PY{c+c1}{t}\PY{c+c1}{i}\PY{c+c1}{o}\PY{c+c1}{n}\PY{c+c1}{ }\PY{c+c1}{2}\PY{c+c1}{ }\PY{c+c1}{s}\PY{c+c1}{t}\PY{c+c1}{r}\PY{c+c1}{i}\PY{c+c1}{c}\PY{c+c1}{t}\PY{c+c1}{l}\PY{c+c1}{y}\PY{c+c1}{ }\PY{c+c1}{d}\PY{c+c1}{o}\PY{c+c1}{m}\PY{c+c1}{i}\PY{c+c1}{n}\PY{c+c1}{a}\PY{c+c1}{n}\PY{c+c1}{t}\PY{c+c1}{ }\PY{c+c1}{f}\PY{c+c1}{o}\PY{c+c1}{r}\PY{c+c1}{ }\PY{c+c1}{b}\PY{c+c1}{o}\PY{c+c1}{t}\PY{c+c1}{h}\PY{c+c1}{ }\PY{c+c1}{p}\PY{c+c1}{l}\PY{c+c1}{a}\PY{c+c1}{y}\PY{c+c1}{e}\PY{c+c1}{r}\PY{c+c1}{s}
\PY{k}{replace} pay1 = \PY{l+m}{3}\PY{k}{ if} game\PY{o}{==}\PY{l+m}{2} \PY{o}{\PYZam{}} player\PY{o}{==}\PY{l+m}{1}
\PY{k}{replace} pay2 = \PY{l+m}{0}\PY{k}{ if} game\PY{o}{==}\PY{l+m}{2} \PY{o}{\PYZam{}} player\PY{o}{==}\PY{l+m}{1}
\PY{k}{replace} pay3 = \PY{l+m}{5}\PY{k}{ if} game\PY{o}{==}\PY{l+m}{2} \PY{o}{\PYZam{}} player\PY{o}{==}\PY{l+m}{1}
\PY{k}{replace} pay4 = \PY{l+m}{1}\PY{k}{ if} game\PY{o}{==}\PY{l+m}{2} \PY{o}{\PYZam{}} player\PY{o}{==}\PY{l+m}{1}

\PY{k}{replace} pay1 = \PY{l+m}{3}\PY{k}{ if} game\PY{o}{==}\PY{l+m}{2} \PY{o}{\PYZam{}} player\PY{o}{==}\PY{l+m}{2}
\PY{k}{replace} pay2 = \PY{l+m}{5}\PY{k}{ if} game\PY{o}{==}\PY{l+m}{2} \PY{o}{\PYZam{}} player\PY{o}{==}\PY{l+m}{2}
\PY{k}{replace} pay3 = \PY{l+m}{0}\PY{k}{ if} game\PY{o}{==}\PY{l+m}{2} \PY{o}{\PYZam{}} player\PY{o}{==}\PY{l+m}{2}
\PY{k}{replace} pay4 = \PY{l+m}{1}\PY{k}{ if} game\PY{o}{==}\PY{l+m}{2} \PY{o}{\PYZam{}} player\PY{o}{==}\PY{l+m}{2}

\PY{c+c1}{*}\PY{c+c1}{ }\PY{c+c1}{G}\PY{c+c1}{a}\PY{c+c1}{m}\PY{c+c1}{e}\PY{c+c1}{ }\PY{c+c1}{3}\PY{c+c1}{:}\PY{c+c1}{ }\PY{c+c1}{m}\PY{c+c1}{a}\PY{c+c1}{t}\PY{c+c1}{c}\PY{c+c1}{h}\PY{c+c1}{i}\PY{c+c1}{n}\PY{c+c1}{g}\PY{c+c1}{ }\PY{c+c1}{p}\PY{c+c1}{e}\PY{c+c1}{n}\PY{c+c1}{n}\PY{c+c1}{i}\PY{c+c1}{e}\PY{c+c1}{s}
\PY{k}{replace} pay1 =  \PY{l+m}{1}\PY{k}{ if} game\PY{o}{==}\PY{l+m}{3} \PY{o}{\PYZam{}} player\PY{o}{==}\PY{l+m}{1}
\PY{k}{replace} pay2 = \PY{o}{\PYZhy{}}\PY{l+m}{1}\PY{k}{ if} game\PY{o}{==}\PY{l+m}{3} \PY{o}{\PYZam{}} player\PY{o}{==}\PY{l+m}{1}
\PY{k}{replace} pay3 = \PY{o}{\PYZhy{}}\PY{l+m}{1}\PY{k}{ if} game\PY{o}{==}\PY{l+m}{3} \PY{o}{\PYZam{}} player\PY{o}{==}\PY{l+m}{1}
\PY{k}{replace} pay4 =  \PY{l+m}{1}\PY{k}{ if} game\PY{o}{==}\PY{l+m}{3} \PY{o}{\PYZam{}} player\PY{o}{==}\PY{l+m}{1}

\PY{k}{replace} pay1 = \PY{o}{\PYZhy{}}\PY{l+m}{1}\PY{k}{ if} game\PY{o}{==}\PY{l+m}{3} \PY{o}{\PYZam{}} player\PY{o}{==}\PY{l+m}{2}
\PY{k}{replace} pay2 =  \PY{l+m}{1}\PY{k}{ if} game\PY{o}{==}\PY{l+m}{3} \PY{o}{\PYZam{}} player\PY{o}{==}\PY{l+m}{2}
\PY{k}{replace} pay3 =  \PY{l+m}{1}\PY{k}{ if} game\PY{o}{==}\PY{l+m}{3} \PY{o}{\PYZam{}} player\PY{o}{==}\PY{l+m}{2}
\PY{k}{replace} pay4 = \PY{o}{\PYZhy{}}\PY{l+m}{1}\PY{k}{ if} game\PY{o}{==}\PY{l+m}{3} \PY{o}{\PYZam{}} player\PY{o}{==}\PY{l+m}{2}

dagamesolve s1\PY{o}{\PYZhy{}}s4, group(game) payoffs(pay1\PY{o}{\PYZhy{}}pay4) equilibria(e)

\PY{k}{list} game player e\PYZus{}count\PY{k}{ if} \PY{o}{!}\PY{n+nf}{missing}(e\PYZus{}count), noobs sep(\PY{l+m}{0})

\PY{k}{assert} e\PYZus{}count\PY{o}{==}\PY{l+m}{3}\PY{k}{ if} game\PY{o}{==}\PY{l+m}{1} \PY{o}{\PYZam{}} player\PY{o}{==}\PY{l+m}{2}
\PY{k}{assert} e\PYZus{}count\PY{o}{==}\PY{l+m}{1}\PY{k}{ if} game\PY{o}{==}\PY{l+m}{2} \PY{o}{\PYZam{}} player\PY{o}{==}\PY{l+m}{2}
\PY{k}{assert} e\PYZus{}count\PY{o}{==}\PY{l+m}{1}\PY{k}{ if} game\PY{o}{==}\PY{l+m}{3} \PY{o}{\PYZam{}} player\PY{o}{==}\PY{l+m}{2}

\PY{k}{display}\PY{k}{ as} result \PY{l+s}{\PYZdq{}}\PY{l+s}{P}\PY{l+s}{A}\PY{l+s}{S}\PY{l+s}{S}\PY{l+s}{:}\PY{l+s}{ }\PY{l+s}{s}\PY{l+s}{e}\PY{l+s}{q}\PY{l+s}{u}\PY{l+s}{e}\PY{l+s}{n}\PY{l+s}{c}\PY{l+s}{e}\PY{l+s}{ }\PY{l+s}{o}\PY{l+s}{f}\PY{l+s}{ }\PY{l+s}{g}\PY{l+s}{a}\PY{l+s}{m}\PY{l+s}{e}\PY{l+s}{s}\PY{l+s}{ }\PY{l+s}{s}\PY{l+s}{o}\PY{l+s}{l}\PY{l+s}{v}\PY{l+s}{e}\PY{l+s}{d}\PY{l+s}{ }\PY{l+s}{c}\PY{l+s}{o}\PY{l+s}{r}\PY{l+s}{r}\PY{l+s}{e}\PY{l+s}{c}\PY{l+s}{t}\PY{l+s}{l}\PY{l+s}{y}\PY{l+s}{\PYZdq{}}
\end{Verbatim}
\end{tcolorbox}

    \begin{Verbatim}[commandchars=\\\{\}]
Number of observations (\_N) was 0, now 6.
(6 missing values generated)
(6 missing values generated)
(6 missing values generated)
(6 missing values generated)
(2 real changes made)
(2 real changes made)
(2 real changes made)
(2 real changes made)
(1 real change made)
(1 real change made)
(1 real change made)
(1 real change made)
(1 real change made)
(1 real change made)
(1 real change made)
(1 real change made)
(1 real change made)
(1 real change made)
(1 real change made)
(1 real change made)
(1 real change made)
(1 real change made)
(1 real change made)
(1 real change made)

  +-------------------------+
  | game   player   e\_count |
  |-------------------------|
  |    1        2         3 |
  |    2        2         1 |
  |    3        2         1 |
  +-------------------------+
PASS: sequence of games solved correctly
    \end{Verbatim}

    \section{Stored results}\label{stored-results}

\texttt{dagamesolve} stores summary information in \texttt{r()} in
addition to the variables it creates in the dataset.

Core stored results include:

\begin{itemize}
\tightlist
\item
  \texttt{r(playeractions)}: number of actions available to each player;
\item
  \texttt{r(totalactions)}: total number of player-actions;
\item
  \texttt{r(method)}: \texttt{complete} or \texttt{newton};
\item
  \texttt{r(games)}: number of games solved;
\item
  \texttt{r(meaneqs)}: mean equilibrium count across games;
\item
  \texttt{r(maxeqs)}: maximum equilibrium count;
\item
  \texttt{r(mineqs)}: minimum equilibrium count;
\item
  \texttt{r(sdeqs)}: standard deviation of equilibrium counts.
\end{itemize}

The complete method also records the interval and Newton settings. The
fast method records the number of random starting points.

Because the equilibrium probability variables remain in the dataset,
they can be reshaped, summarized, merged to market-level data, or used
in subsequent programs without leaving Stata.

    \begin{tcolorbox}[breakable, size=fbox, boxrule=1pt, pad at break*=1mm,colback=cellbackground, colframe=cellborder]
\prompt{In}{incolor}{102}{\boxspacing}
\begin{Verbatim}[commandchars=\\\{\}]
\PY{k}{return}\PY{k}{ list}
\end{Verbatim}
\end{tcolorbox}

    \begin{Verbatim}[commandchars=\\\{\}]

scalars:
              r(sdeqs) =  1.154700538379252
             r(mineqs) =  1
             r(maxeqs) =  3
            r(meaneqs) =  1.666666666666667
              r(games) =  3
          r(tolerance) =  1.00000000000e-12
             r(digits) =  1.00000000000e-08
              r(mnits) =  100
             r(coltol) =  .0001
            r(mbwidth) =  .0001
       r(totalactions) =  4

macros:
             r(method) : "complete"

matrices:
      r(playeractions) :  1 x 2
    \end{Verbatim}

    \section{Computational
considerations}\label{computational-considerations}

The computational burden grows quickly for two related reasons.

First, the number of pure profiles is

\[
\prod_i K_i.
\]

Second, mixed-equilibrium search requires consideration of candidate
supports. For each surviving support, the solver constructs and solves a
nonlinear system. Adding players or actions therefore affects both the
size of individual systems and the number of systems that may have to be
considered.

Interval methods are intentionally more demanding than point-based
Newton searches. Their advantage is that boxes can be ruled out over
entire regions of the strategy space rather than only explored from
selected starting points.

Several practical implications follow:

\begin{itemize}
\tightlist
\item
  use \texttt{pureonly} when mixed equilibria are irrelevant;
\item
  use \texttt{fast} when exploratory speed matters more than systematic
  coverage;
\item
  use the default method when the goal is a systematic search for pure
  and mixed equilibria;
\item
  exploit sequences of games when many games have the same action
  structure;
\item
  inspect difficult, degenerate, or nearly duplicate equilibria
  carefully;
\item
  report numerical settings in applications where equilibrium
  multiplicity is substantively important.
\end{itemize}

    \section{What this makes possible for empirical
work}\label{what-this-makes-possible-for-empirical-work}

The original motivation for these routines was closely connected to
empirical models of strategic interaction.

Suppose market \(m\) is associated with a game whose payoff primitives
depend on observed covariates \(X_m\) and parameters \(\theta\):

\[
\Pi_m=\Pi(X_m,\theta).
\]

For a trial value of \(\theta\), one can:

\begin{enumerate}
\def\labelenumi{\arabic{enumi}.}
\tightlist
\item
  construct the payoff matrix for every market;
\item
  solve the sequence of games;
\item
  obtain the equilibrium correspondence market by market;
\item
  compare predicted equilibrium behavior with observed outcomes;
\item
  embed the calculation in an estimation, simulation, or counterfactual
  procedure.
\end{enumerate}

The computational problem is therefore not incidental. When a structural
model permits multiple equilibria, the ability to obtain the equilibrium
correspondence repeatedly can be an essential piece of the empirical
method.

Bajari, Hong, and Ryan (2010) provide an important example of the
broader connection between discrete games and econometric analysis.

    \section{Conclusions}\label{conclusions}

\texttt{dagamesolve} is designed to make finite-game equilibrium
computation available inside a data-analysis environment.

The easy part of a finite game is often the part one sees in a textbook:
enumerate a few pure profiles and check unilateral deviations. The
difficult part is what happens once mixed and partially mixed equilibria
are admitted. The problem becomes one of support enumeration and
nonlinear equation solving.

The package addresses that problem by combining game-specific reductions
with general numerical tools:

\begin{itemize}
\tightlist
\item
  dominance reduction;
\item
  direct pure-strategy checks;
\item
  support enumeration;
\item
  interval contraction and subdivision;
\item
  Newton refinement;
\item
  candidate verification and consolidation.
\end{itemize}

The ability to apply the same machinery to a sequence of games is
particularly useful for applications in which payoffs vary across
markets, observations, or counterfactual parameter values.

The present notebook is intended to serve simultaneously as
documentation, a reproducible collection of examples, and a technical
description of the solver.

    \section*{References}\label{references}
\addcontentsline{toc}{section}{References}

Aradillas-López, Andrés. 2020. ``The Econometrics of Static Games.''
\emph{Annual Review of Economics} 12: 135--165.

Bajari, Patrick, Han Hong, John Krainer, and Denis Nekipelov. 2010.
``Estimating Static Models of Strategic Interactions.'' \emph{Journal of
Business \& Economic Statistics} 28(4): 469--482.

Bajari, Patrick, Han Hong, and Stephen P. Ryan. 2010. ``Identification
and Estimation of a Discrete Game of Complete Information.''
\emph{Econometrica} 78(5): 1529--1568.

Berg, Johannes, and Andrew McLennan. 2005. ``Asymptotic Expected Number
of Nash Equilibria of Two-Player Normal Form Games.'' \emph{Games and
Economic Behavior} 51(2): 264--295.

Berry, Steven T. 1992. ``Estimation of a Model of Entry in the Airline
Industry.'' \emph{Econometrica} 60(4): 889--917.

Bresnahan, Timothy F., and Peter C. Reiss. 1990. ``Entry in Monopoly
Market.'' \emph{Review of Economic Studies} 57(4): 531--553.

Bresnahan, Timothy F., and Peter C. Reiss. 1991. ``Entry and Competition
in Concentrated Markets.'' \emph{Journal of Political Economy} 99(5):
977--1009.

Ciliberto, Federico, and Elie Tamer. 2009. ``Market Structure and
Multiple Equilibria in Airline Markets.'' \emph{Econometrica} 77(6):
1791--1828.

Conitzer, Vincent, and Tuomas Sandholm. 2008. ``New Complexity Results
about Nash Equilibria.'' \emph{Games and Economic Behavior} 63(2):
621--641.

Daskalakis, Constantinos, Paul W. Goldberg, and Christos H.
Papadimitriou. 2009. ``The Complexity of Computing a Nash Equilibrium.''
\emph{SIAM Journal on Computing} 39(1): 195--259.

Etessami, Kousha, and Mihalis Yannakakis. 2010. ``On the Complexity of
Nash Equilibria and Other Fixed Points.'' \emph{SIAM Journal on
Computing} 39(6): 2531--2597.

Gilboa, Itzhak, and Eitan Zemel. 1989. ``Nash and Correlated Equilibria:
Some Complexity Considerations.'' \emph{Games and Economic Behavior}
1(1): 80--93.

Grieco, Paul L. E. 2014. ``Discrete Games with Flexible Information
Structures: An Application to Local Grocery Markets.'' \emph{RAND
Journal of Economics} 45(2): 303--340.

Jaulin, Luc, Michel Kieffer, Olivier Didrit, and Eric Walter. 2001.
\emph{Applied Interval Analysis}. London: Springer.

Jia, Panle. 2008. ``What Happens When Wal-Mart Comes to Town: An
Empirical Analysis of the Discount Retailing Industry.''
\emph{Econometrica} 76(6): 1263--1316.

Mazzeo, Michael J. 2002. ``Product Choice and Oligopoly Market
Structure.'' \emph{RAND Journal of Economics} 33(2): 221--242.

McKelvey, Richard D. 1998. ``A Liapunov Function for Nash Equilibria.''
\emph{Social Science Working Paper 953}, California Institute of
Technology.

McKelvey, Richard D., and Andrew McLennan. 1996. ``Computation of
Equilibria in Finite Games.'' In \emph{Handbook of Computational
Economics}, Vol. 1, edited by Hans M. Amman, David A. Kendrick, and John
Rust, 87--142. Amsterdam: North-Holland.

McKelvey, Richard D., Andrew M. McLennan, and Theodore L. Turocy. 2014.
\emph{Gambit: Software Tools for Game Theory, Version 14.1.1.} The
Gambit Project.

McLennan, Andrew. 2005. ``The Expected Number of Nash Equilibria of a
Normal Form Game.'' \emph{Econometrica} 73(1): 141--174.

de Paula, Áureo. 2013. ``Econometric Analysis of Games with Multiple
Equilibria.'' \emph{Annual Review of Economics} 5: 107--131.

Seim, Katja. 2006. ``An Empirical Model of Firm Entry with Endogenous
Product-Type Choices.'' \emph{RAND Journal of Economics} 37(3):
619--640.

Sweeting, Andrew. 2009. ``The Strategic Timing Incentives of Commercial
Radio Stations: An Empirical Analysis Using Multiple Equilibria.''
\emph{RAND Journal of Economics} 40(4): 710--742.

Wilson, Robert. 1971. ``Computing Equilibria of N-Person Games.''
\emph{SIAM Journal on Applied Mathematics} 21(1): 80--87.

    \section*{Appendix A: Installation and software
environment}\label{appendix-a-installation-and-software-environment}
\addcontentsline{toc}{section}{Appendix A: Installation and software
environment}

The executable examples require Stata, a Jupyter installation with the
\textbf{Stata (nbstata)} kernel, and the Stata/Mata packages
\texttt{dagamesolve}, \texttt{dagamestrats}, \texttt{moremata},
\texttt{int\_utils}, \texttt{rowmat\_utils}, and \texttt{intsolver}.
\texttt{dagamestrats} is supplied with \texttt{dagamesolve}. The
\texttt{dependency.do} file in the DaGameSolve repository also documents
installation of the dependencies.

\textbf{Installation.} If these packages are not already available on
Stata's ado-path, run the following commands in Stata \textbf{before}
running the examples. They are shown as a code listing rather than an
automatically executing cell: this prevents Jupyter's \emph{Run All}
command from unnecessarily reinstalling packages (the \texttt{replace}
option can overwrite an existing installation).

\begin{Shaded}
\begin{Highlighting}[]
\KeywordTok{ssc}\NormalTok{ install moremata, }\KeywordTok{replace}

\NormalTok{net install int\_utils, }\CommentTok{///}
\NormalTok{    from(}\StringTok{"https://raw.githubusercontent.com/mbaker21231/int\_utils/main/"}\NormalTok{) }\KeywordTok{replace}

\NormalTok{net install rowmat\_utils, }\CommentTok{///}
\NormalTok{    from(}\StringTok{"https://raw.githubusercontent.com/mbaker21231/rowmat\_utils/main/"}\NormalTok{) }\KeywordTok{replace}

\NormalTok{net install intsolver, }\CommentTok{///}
\NormalTok{    from(}\StringTok{"https://raw.githubusercontent.com/mbaker21231/IntSolver/main/"}\NormalTok{) }\KeywordTok{replace}

\NormalTok{* DaGameSolve development }\KeywordTok{version}\NormalTok{ (profiling branch)}
\NormalTok{net install dagamesolve, }\CommentTok{///}
\NormalTok{    from(}\StringTok{"https://raw.githubusercontent.com/mbaker21231/DaGameSolve/profiling/"}\NormalTok{) }\KeywordTok{replace}

\NormalTok{* Refresh Mata\textquotesingle{}s library }\FunctionTok{index}
\KeywordTok{mata}\NormalTok{: }\KeywordTok{mata}\NormalTok{ mlib }\FunctionTok{index}
\end{Highlighting}
\end{Shaded}

The DaGameSolve URL above intentionally selects the
\textbf{\texttt{profiling}} branch, matching the installation used while
preparing this notebook; it is not necessarily the released version.

\textbf{Environment and package checks.} The following cell reports the
Stata version and finds the installed implementations of the two main
commands. It can be rerun after installing or updating packages. Using
\texttt{c(stata\_version)} (rather than \texttt{about}) avoids embedding
Stata license information in a notebook that may later be shared or
published.

    \begin{tcolorbox}[breakable, size=fbox, boxrule=1pt, pad at break*=1mm,colback=cellbackground, colframe=cellborder]
\prompt{In}{incolor}{108}{\boxspacing}
\begin{Verbatim}[commandchars=\\\{\}]
\PY{k}{display}\PY{k}{ as} text \PY{l+s}{\PYZdq{}}\PY{l+s}{S}\PY{l+s}{t}\PY{l+s}{a}\PY{l+s}{t}\PY{l+s}{a}\PY{l+s}{ }\PY{l+s}{v}\PY{l+s}{e}\PY{l+s}{r}\PY{l+s}{s}\PY{l+s}{i}\PY{l+s}{o}\PY{l+s}{n}\PY{l+s}{:}\PY{l+s}{ }\PY{l+s}{\PYZdq{}} \PY{n+nf}{c}(stata\PYZus{}version)
\PY{k}{which} dagamesolve
\PY{k}{which} dagamestrats
\end{Verbatim}
\end{tcolorbox}

    \begin{Verbatim}[commandchars=\\\{\}]
Stata version: 19.5
.\textbackslash{}dagamesolve.ado
*! dagsolve v1.0.1 MJBaker 21july2014
*! dagsolve v1.0.2 MJBaker 18sep2014
*! dagamesolve v1.0.3 MJBaker 8may2023
.\textbackslash{}dagamestrats.ado
*! dagstrats v1.0.0 MJBaker 17July2014
*! dagamestrats v2.0.0 MJBaker 30May2023
    \end{Verbatim}

    \begin{tcolorbox}[breakable, size=fbox, boxrule=1pt, pad at break*=1mm,colback=cellbackground, colframe=cellborder]
\prompt{In}{incolor}{ }{\boxspacing}
\begin{Verbatim}[commandchars=\\\{\}]

\end{Verbatim}
\end{tcolorbox}


    % Add a bibliography block to the postdoc
    
    
    
\end{document}
